% Auto-generated by scripts/render_provenance.py -- do not edit by hand.
% Place % Auto-generated by scripts/render_provenance.py -- do not edit by hand.
% Place % Auto-generated by scripts/render_provenance.py -- do not edit by hand.
% Place % Auto-generated by scripts/render_provenance.py -- do not edit by hand.
% Place \input{data/provenance_tables} inside supplementary.tex.
% (supplementary.tex preamble already loads \usepackage{float}.)

\begin{table}[H]
    \caption{Source classification: Piezo1.}
    \label{tab:prov_piezo1}
    \footnotesize
    \centering
    \begin{tabular}{@{}llp{6.5cm}@{}}
        \toprule
        Parameter & Source & Reference / notes \\
        \midrule
        \texttt{Ca\_permeability\_fraction} & Literature & Coste et al. 2010 (P\_Ca / P\_Na \textasciitilde{}5) --- \textasciitilde{}10\% of cation current is Ca\textsuperscript{2}\textsuperscript{+} in physiological gradients \\
        \texttt{E\_rev} & Literature & Coste et al. 2010 (non-selective cation, \textasciitilde{}0 mV) \\
        \texttt{T\_half} & Literature & Coste et al. 2010, Science 330:55-60; Lewis \& Grandl 2015 --- Mid-point of activation across multiple cell types \\
        \texttt{delta\_A} & Literature & Coste et al. 2010 (Boltzmann fit to single-channel records) --- Effective gating area change from $\Delta$A$\cdot$T fit \\
        \texttt{g\_max} & Literature & Coste et al. 2010 (single-channel slope) --- Single-channel conductance \textasciitilde{}25 pS in physiological gradients \\
        \texttt{tau\_act} & Literature & Coste et al. 2010 Fig. 2 (activation rise $\tau$ \textasciitilde{}1-2 ms) --- Phase-5 fixed-tau replacement for the bell-shaped $\tau$\_w(T) \\
        \texttt{tau\_inact} & Literature & Coste et al. 2010 Fig. 2 (inactivation $\tau$ 10-15 ms) \\
        \texttt{g\_density} & Assumption & --- --- Sparse expression assumption; user typically overrides per cell type \\
        \texttt{inact\_T\_half} & Assumption & --- --- Set above activation T\_half; not directly fit in Coste 2010 \\
        \texttt{inact\_delta\_A} & Assumption & --- --- Comparable steepness to activation; placeholder pending dedicated fit \\
        \texttt{tau\_0} & Assumption & --- --- Single time constant subsuming activation kinetics; superseded by tau\_act when set \\
        \bottomrule
    \end{tabular}
\end{table}

\begin{table}[H]
    \caption{Source classification: TRAAK.}
    \label{tab:prov_traak}
    \footnotesize
    \centering
    \begin{tabular}{@{}llp{6.5cm}@{}}
        \toprule
        Parameter & Source & Reference / notes \\
        \midrule
        \texttt{E\_rev} & Literature & K\textsuperscript{+} Nernst at physiological [K] --- -90 mV typical \\
        \texttt{T\_half} & Literature & Brohawn et al. 2014, PNAS 111:3614-3619 \\
        \texttt{delta\_A} & Literature & Brohawn et al. 2014 \\
        \texttt{g\_max} & Literature & Brohawn et al. 2014 \\
        \texttt{tau\_act} & Literature & Brohawn et al. 2014 (single-channel mechanical activation rise) --- Phase-5 fixed-tau replacement for the bell-shaped $\tau$\_w(T) \\
        \texttt{g\_density} & Assumption & Estimated from neural expression levels; user-tunable \\
        \texttt{tau\_0} & Assumption & Legacy bell-shape baseline; superseded by tau\_act on the library preset \\
        \texttt{Ca\_permeability\_fraction} & Derived & K2P channels are K\textsuperscript{+}-selective; no Ca\textsuperscript{2}\textsuperscript{+} flux \\
        \texttt{inact\_T\_half} & Derived & Not applicable; tau\_inact is None \\
        \texttt{inact\_delta\_A} & Derived & Not applicable; tau\_inact is None \\
        \texttt{tau\_inact} & Derived & K2P channels don't inactivate; field defaults to None (gate held at 1) \\
        \bottomrule
    \end{tabular}
\end{table}

\begin{table}[H]
    \caption{Source classification: TREK-1.}
    \label{tab:prov_trek_1}
    \footnotesize
    \centering
    \begin{tabular}{@{}llp{6.5cm}@{}}
        \toprule
        Parameter & Source & Reference / notes \\
        \midrule
        \texttt{E\_rev} & Literature & K\textsuperscript{+} Nernst \\
        \texttt{T\_half} & Literature & Honor\'{e} 2007, Nat Rev Neurosci 8:251-261 \\
        \texttt{delta\_A} & Literature & Honor\'{e} 2007 \\
        \texttt{g\_max} & Literature & Honor\'{e} 2007 \\
        \texttt{tau\_act} & Literature & Honor\'{e} 2007 (mid-range of 3-10 ms activation $\tau$ reported in review) --- Phase-5 fixed-tau replacement for the bell-shaped $\tau$\_w(T) \\
        \texttt{g\_density} & Assumption & Cortical pyramidal neuron expression range \\
        \texttt{tau\_0} & Assumption & Legacy bell-shape baseline; superseded by tau\_act on the library preset \\
        \texttt{Ca\_permeability\_fraction} & Derived & K2P channels are K\textsuperscript{+}-selective; no Ca\textsuperscript{2}\textsuperscript{+} flux \\
        \texttt{inact\_T\_half} & Derived & Not applicable; tau\_inact is None \\
        \texttt{inact\_delta\_A} & Derived & Not applicable; tau\_inact is None \\
        \texttt{tau\_inact} & Derived & K2P channels don't inactivate; field defaults to None \\
        \bottomrule
    \end{tabular}
\end{table}

\begin{table}[H]
    \caption{Source classification: TREK-2.}
    \label{tab:prov_trek_2}
    \footnotesize
    \centering
    \begin{tabular}{@{}llp{6.5cm}@{}}
        \toprule
        Parameter & Source & Reference / notes \\
        \midrule
        \texttt{E\_rev} & Literature & K\textsuperscript{+} Nernst \\
        \texttt{T\_half} & Literature & Kang \& Kim 2024 review \\
        \texttt{delta\_A} & Literature & Kang \& Kim 2024 \\
        \texttt{g\_max} & Literature & Kang \& Kim 2024 \\
        \texttt{tau\_act} & Literature & Kang \& Kim 2024 (TREK-2 kinetics; similar to TREK-1) --- Phase-5 fixed-tau replacement for the bell-shaped $\tau$\_w(T) \\
        \texttt{g\_density} & Assumption & Cortical pyramidal neuron expression range \\
        \texttt{tau\_0} & Assumption & Legacy bell-shape baseline; superseded by tau\_act on the library preset \\
        \texttt{Ca\_permeability\_fraction} & Derived & K2P channels are K\textsuperscript{+}-selective; no Ca\textsuperscript{2}\textsuperscript{+} flux \\
        \texttt{inact\_T\_half} & Derived & Not applicable; tau\_inact is None \\
        \texttt{inact\_delta\_A} & Derived & Not applicable; tau\_inact is None \\
        \texttt{tau\_inact} & Derived & K2P channels don't inactivate; field defaults to None \\
        \bottomrule
    \end{tabular}
\end{table}

\begin{table}[H]
    \caption{Source classification: TRPV1.}
    \label{tab:prov_trpv1}
    \footnotesize
    \centering
    \begin{tabular}{@{}llp{6.5cm}@{}}
        \toprule
        Parameter & Source & Reference / notes \\
        \midrule
        \texttt{E\_rev} & Literature & Slightly Ca\textsuperscript{2}\textsuperscript{+}-preferring non-selective cation, \textasciitilde{}10 mV \\
        \texttt{T\_half} & Literature & Caterina et al. 1997, 2000 (heat activation threshold) \\
        \texttt{k} & Calibrated & Boltzmann steepness chosen so Po(47)/Po(37) > 20 matching the literature Q10 $\geq$ 20 floor (Caterina 1997). \\
        \texttt{g\_density} & Assumption & Conservative mid-range density when expressed \\
        \bottomrule
    \end{tabular}
\end{table}

\begin{table}[H]
    \caption{Source classification: TRPV4.}
    \label{tab:prov_trpv4}
    \footnotesize
    \centering
    \begin{tabular}{@{}llp{6.5cm}@{}}
        \toprule
        Parameter & Source & Reference / notes \\
        \midrule
        \texttt{E\_rev} & Literature & Non-selective cation, \textasciitilde{}0 mV \\
        \texttt{T\_half} & Literature & Watanabe et al. 2002 (warm-activation threshold) --- Within the 27-35$^\circ$C literature range \\
        \texttt{k} & Calibrated & Q10 $\approx$ 10 around threshold, less steep than TRPV1 \\
        \texttt{g\_density} & Assumption & Conservative mid-range density when expressed \\
        \bottomrule
    \end{tabular}
\end{table}

\begin{table}[H]
    \caption{Source classification: TRPA1.}
    \label{tab:prov_trpa1}
    \footnotesize
    \centering
    \begin{tabular}{@{}llp{6.5cm}@{}}
        \toprule
        Parameter & Source & Reference / notes \\
        \midrule
        \texttt{Ca\_permeability\_fraction} & Literature & Karashima et al. 2010; Wang et al. 2008 (P\_Ca/P\_Na $\approx$ 5--8) --- \textasciitilde{}15--20\% of TRPA1 cation current is Ca\textsuperscript{2}\textsuperscript{+} at physiological gradients. \\
        \texttt{E\_rev} & Literature & Karashima et al. 2010 (non-selective cation, E\_rev $\approx$ 0 mV) \\
        \texttt{g\_max} & Literature & Karashima et al. 2010; Wang et al. 2008 (TRPA1 single-channel \textasciitilde{}tens of pS) \\
        \texttt{g\_density} & Assumption & Astrocytic TRPA1 surface density; user overrides per preparation. \\
        \texttt{tau\_0} & Assumption & Legacy bell-shape baseline; superseded by tau\_act. \\
        \texttt{T\_half} & Estimate & Oh et al. 2019, Curr Biol 29:3386 (US activates astrocytic TRPA1) --- Half-activation tension for TRPA1 mechanogating is not directly quantified; set between Piezo1 (2.7) and TRAAK (5.0) mN/m as a best guess pending a dedicated tension-clamp fit. \\
        \texttt{delta\_A} & Estimate & Gating-area sensitivity guessed comparable to other MS channels (\textasciitilde{}15 nm\textsuperscript{2}). \\
        \texttt{tau\_act} & Estimate & Channel gating is fast (ms); the slow relay signature lives in the downstream Ca/glutamate dynamics, not in TRPA1 gating. \\
        \bottomrule
    \end{tabular}
\end{table}

\begin{table}[H]
    \caption{Source classification: strain\_to\_tension.}
    \label{tab:prov_strain_to_tension}
    \footnotesize
    \centering
    \begin{tabular}{@{}llp{6.5cm}@{}}
        \toprule
        Parameter & Source & Reference / notes \\
        \midrule
        \texttt{K\_A} & Assumption & Lipid-bilayer area-expansion modulus, range 0.1-0.5 N/m --- Default 0.25 N/m is the centre of the published range (Rawicz 2000, Sukharev 2012); the actual coupling from tissue-scale strain to per-channel membrane tension subsumes geometric concentration, cytoskeletal mechanics, and membrane curvature effects. User must calibrate for the specific tissue/cell-type/exposure under study. \\
        \texttt{coupling\_factor} & Assumption & Dimensionless multi-scale amplification $\alpha$ relating far-field tissue strain to local membrane area change. Code default 1.0; manuscript-canonical $\alpha$=1000 (cortical-pyramidal-calibrated, 5$\times$ the iso-potential/hippocampal literature $\alpha$$\approx$200). Calibrate per tissue/cell-type/exposure. \\
        \bottomrule
    \end{tabular}
\end{table}

\begin{table}[H]
    \caption{Source classification: NICEModel.}
    \label{tab:prov_nicemodel}
    \footnotesize
    \centering
    \begin{tabular}{@{}llp{6.5cm}@{}}
        \toprule
        Parameter & Source & Reference / notes \\
        \midrule
        \texttt{C\_m0} & Literature & Standard neuronal specific capacitance, 1 $\mu$F/cm\textsuperscript{2} \\
        \texttt{delta\_0} & Literature & Lipid bilayer half-thickness (Plaksin 2014) \\
        \texttt{f\_carrier} & Literature & Typical TUS carrier (Legon, Tyler) --- 500 kHz default; user supplies their protocol value \\
        \texttt{f\_resonant} & Literature & Plaksin et al. 2014 (leaflet resonance \textasciitilde{}5 MHz) \\
        \texttt{lj\_strength} & Calibrated & Dimensionless prefactor for the LJ-like wall on Z<0; calibrated so the wall contribution stays <5\% of linear restoring at |Z| < 0.1$\cdot$$\delta$\_0 \\
        \texttt{m\_eff} & Calibrated & SONIC operating range, Lemaire et al. 2019 --- Calibrated so |Z|(100 kPa) $\approx$ 0.6 nm $\to$ $\langle$$\Delta$C\_m$\rangle$/C\_m0 $\approx$ 5\% \\
        \texttt{st\_strength} & Calibrated & Dimensionless prefactor for the surface-tension wall on Z>0; default 1.0 doubles asymptotic restoring on the expansion side, extending the quantitative validity range to \textasciitilde{}1 MPa (Plaksin \& Krasovitski 2014 Fig. 2). \\
        \texttt{damping\_ratio} & Assumption & $\zeta$ = 0.1 (Q $\approx$ 5); Plaksin-consistent moderate damping \\
        \bottomrule
    \end{tabular}
\end{table}

\begin{table}[H]
    \caption{Source classification: SolitonModel.}
    \label{tab:prov_solitonmodel}
    \footnotesize
    \centering
    \begin{tabular}{@{}llp{6.5cm}@{}}
        \toprule
        Parameter & Source & Reference / notes \\
        \midrule
        \texttt{T\_m} & Literature & Heimburg 2007, Thermal Biophysics of Membranes; Heimburg \& Jackson 2005, PNAS 102:9790 --- Membrane melting transition. Biological membranes melt \textasciitilde{}10-15$^\circ$C below body temperature; default 30$^\circ$C is representative (lung-surfactant T\_m $\approx$ 30$^\circ$C, the system Heimburg \& Jackson fit). Must be < 37$^\circ$C for the heating-reduces-excitability sign. \\
        \texttt{c} & Literature & Soft-tissue speed of sound; matches repo FDTD C0 = 1540 m/s --- $\rho$c\textsuperscript{2} is the adiabatic bulk modulus for the strain conversion. \\
        \texttt{rho} & Literature & Soft-tissue/water density; matches repo FDTD RHO0 = 1000 kg/m\textsuperscript{3} --- Used in the bulk-modulus strain $\varepsilon$ = P/($\rho$c\textsuperscript{2}). \\
        \texttt{peak\_shape} & Assumption & Gaussian (default) vs Lorentzian shape for the normalised c\_p peak. \\
        \texttt{c\_p\_max} & Estimate & Heimburg 2007 ($\kappa$\_A\^{}T $\propto$ $\Delta$c\_p proportionality) --- Dimensionless peak-to-baseline excess-compressibility ratio encoding Heimburg's c\_p $\leftrightarrow$ compressibility divergence near T\_m. Tuned amplitude (default 10$\times$). \\
        \texttt{sigma\_T} & Estimate & Heimburg 2007 (broadened multicomponent-membrane transition) --- Half-width of the c\_p / compressibility peak, a few $^\circ$C. Biological transitions are broad; specific value tuned. \\
        \bottomrule
    \end{tabular}
\end{table}

\begin{table}[H]
    \caption{Source classification: CalciumDynamics.}
    \label{tab:prov_calciumdynamics}
    \footnotesize
    \centering
    \begin{tabular}{@{}llp{6.5cm}@{}}
        \toprule
        Parameter & Source & Reference / notes \\
        \midrule
        \texttt{Ca\_rest} & Literature & Cortical pyramidal resting [Ca\textsuperscript{2}\textsuperscript{+}]\_i, \textasciitilde{}50 nM \\
        \texttt{tau\_decay} & Literature & Mainen \& Sejnowski 1996; Borg-Graham 1999 --- Single-pool decay $\tau$ in cortical models, range 30-80 ms \\
        \texttt{flux\_factor} & Calibrated & Lumped sub-membrane shell volume $\times$ buffering fraction $\times$ 1/(z$\cdot$F). Default 0.5 lies within published current-to-concentration coupling factors used in cortical Ca\textsuperscript{2}\textsuperscript{+}-pool models. \\
        \bottomrule
    \end{tabular}
\end{table}

\begin{table}[H]
    \caption{Source classification: KCaChannel.}
    \label{tab:prov_kcachannel}
    \footnotesize
    \centering
    \begin{tabular}{@{}llp{6.5cm}@{}}
        \toprule
        Parameter & Source & Reference / notes \\
        \midrule
        \texttt{Ca\_half} & Literature & K\"{o}hler et al. 1996 (SK channel) \\
        \texttt{E\_K} & Literature & K\textsuperscript{+} Nernst, matches Pospischil RS \\
        \texttt{n\_hill} & Literature & Hirschberg et al. 1998 \\
        \texttt{g\_density} & Assumption & Cortical pyramidal SK density range, 0.1--1.0 mS/cm\textsuperscript{2} (Sah 1996) \\
        \bottomrule
    \end{tabular}
\end{table}

\begin{table}[H]
    \caption{Source classification: AstrocyteRelay.}
    \label{tab:prov_astrocyterelay}
    \footnotesize
    \centering
    \begin{tabular}{@{}llp{6.5cm}@{}}
        \toprule
        Parameter & Source & Reference / notes \\
        \midrule
        \texttt{Ca\_rest} & Literature & Bazargani \& Attwell 2016 (resting astrocytic [Ca\textsuperscript{2}\textsuperscript{+}] \textasciitilde{}50--100 nM) \\
        \texttt{V\_astro} & Literature & Bazargani \& Attwell 2016, Nat Neurosci (astrocyte V\_rest $\approx$ -80 mV) --- Sets the Ca\textsuperscript{2}\textsuperscript{+} driving force through TRPA1 (non-spiking cell). \\
        \texttt{tau\_decay} & Literature & Bazargani \& Attwell 2016; Oh et al. 2019 (astrocytic Ca\textsuperscript{2}\textsuperscript{+} transients decay over 100s of ms) --- Slow relay: hundreds of ms, not per-carrier-cycle. \\
        \texttt{flux\_factor} & Calibrated & Chosen so the astrocytic Ca\textsuperscript{2}\textsuperscript{+} pool peaks \textasciitilde{}1 $\mu$M at full US drive (physiological gliotransmission regime), given TRPA1 g\_density, Ca permeability, and the -80 mV driving force. \\
        \texttt{glu\_Ca\_half} & Estimate & Parpura \& Haydon 2000 (Ca-dependent gliotransmission threshold) --- Half-release [Ca\textsuperscript{2}\textsuperscript{+}] a few hundred nM; tuned for gliotransmission onset. \\
        \texttt{glu\_max} & Estimate & Peak cleft glutamate scale [$\mu$M] from gliotransmission; order-of-magnitude. \\
        \texttt{glu\_n\_hill} & Estimate & Cooperative Ca\textsuperscript{2}\textsuperscript{+}-triggered release; Hill coefficient \textasciitilde{}2--3 (vesicular sensor). \\
        \bottomrule
    \end{tabular}
\end{table}

\begin{table}[H]
    \caption{Source classification: NMDAReceptor.}
    \label{tab:prov_nmdareceptor}
    \footnotesize
    \centering
    \begin{tabular}{@{}llp{6.5cm}@{}}
        \toprule
        Parameter & Source & Reference / notes \\
        \midrule
        \texttt{E\_rev} & Literature & NMDA non-selective cation, E\_rev $\approx$ 0 mV \\
        \texttt{K\_glu} & Literature & Patneau \& Mayer 1990 (NMDA glutamate EC50 \textasciitilde{}1--3 $\mu$M) \\
        \texttt{Mg} & Literature & Jahr \& Stevens 1990 (physiological extracellular [Mg\textsuperscript{2}\textsuperscript{+}] $\approx$ 1 mM) \\
        \texttt{mg\_block\_v\_scale} & Literature & Jahr \& Stevens 1990 (voltage-dependent Mg\textsuperscript{2}\textsuperscript{+} block: 16.13 mV, 3.57 mM) --- g(V) $\propto$ 1/(1 + [Mg]$\cdot$exp(-V/16.13)/3.57). \\
        \texttt{alpha\_NMDA} & Calibrated & Maneshi et al. 2017, Sci Rep (membrane stretch activates NMDARs); Singh et al. 2012 --- Max fractional stretch enhancement of NMDA current. Default 1.0 $\to$ current roughly doubles at saturating stretch, matching Maneshi 2017 (stretch $\approx$ doubled NMDA current). f\_stretch($\infty$) = 1 + alpha\_NMDA. \\
        \texttt{g\_max} & Assumption & Post-synaptic NMDA conductance density; user overrides per cell type. \\
        \texttt{T\_half\_stretch} & Estimate & Maneshi et al. 2017 --- Half-enhancement membrane tension [mN/m]. Mapped to the MS-channel tension range pending a dedicated stretch-clamp NMDA fit. \\
        \texttt{k\_stretch} & Estimate & Boltzmann slope [mN/m] of the stretch sigmoid; comparable steepness to MS channels. \\
        \bottomrule
    \end{tabular}
\end{table}

\begin{table}[H]
    \caption{Source classification: MechanoNav.}
    \label{tab:prov_mechanonav}
    \footnotesize
    \centering
    \begin{tabular}{@{}llp{6.5cm}@{}}
        \toprule
        Parameter & Source & Reference / notes \\
        \midrule
        \texttt{k\_a} & Literature & Hodgkin-Huxley-type Nav activation Boltzmann slope --- Activation slope factor [mV]. \\
        \texttt{v\_half0} & Literature & Representative mammalian Nav activation midpoint (Hodgkin-Huxley-type) --- Baseline activation half-voltage [mV] for the scalar gain. \\
        \texttt{k\_mechano} & Calibrated & Kubanek et al. 2016, Sci Rep (US shifts Nav/Cav voltage dependence) --- Hyperpolarising activation-shift sensitivity [mV per (mN/m)]. K\_MECHANO\_NAV = 0.3 is calibrated so a few-mN/m tension yields a \textasciitilde{}mV-scale activation shift and a tens-of-percent Na\textsuperscript{+} current change, the order Kubanek report (not a directly tabulated value). The neuron field k\_mechano\_nav defaults to 0 (pathway off) to preserve the baseline neuron; set it to K\_MECHANO\_NAV to enable. \\
        \texttt{v\_eval} & Assumption & Near-threshold reference voltage [mV] at which nav\_current\_gain reads the open fraction. \\
        \bottomrule
    \end{tabular}
\end{table}

\begin{table}[H]
    \caption{Source classification: MechanoSynapse.}
    \label{tab:prov_mechanosynapse}
    \footnotesize
    \centering
    \begin{tabular}{@{}llp{6.5cm}@{}}
        \toprule
        Parameter & Source & Reference / notes \\
        \midrule
        \texttt{E\_ampa} & Literature & AMPA non-selective cation reversal, \textasciitilde{}0 mV \\
        \texttt{K\_ampa} & Literature & AMPA is low-affinity (glutamate EC50 \textasciitilde{} hundreds of $\mu$M) --- AMPA glutamate half-saturation [$\mu$M]. \\
        \texttt{Pr\_baseline} & Literature & Tyler 2008; central-synapse release probability range 0.1--0.3 --- Resting vesicle release probability at zero tension. \\
        \texttt{tau\_glu} & Literature & Cleft glutamate clearance / uptake time constant (few ms) --- First-order glutamate pool decay [ms]. \\
        \texttt{g\_ampa} & Calibrated & Post-synaptic AMPA conductance density [mS/cm\textsuperscript{2}]; calibrated with glu\_scale. \\
        \texttt{glu\_scale} & Calibrated & Sustained cleft glutamate [$\mu$M] at full release; calibrated with g\_ampa so the resting (baseline-Pr) state is sub-threshold and US-driven release brings the post neuron to firing. \\
        \texttt{Pr\_max} & Estimate & Tyler 2012 (mechanical perturbation raises release probability) --- Saturating release probability at high presynaptic tension. \\
        \texttt{T\_half\_pre} & Estimate & Tyler 2012 --- Half-modulation presynaptic tension [mN/m]; in the MS-channel range. \\
        \texttt{k\_pre} & Estimate & Boltzmann slope [mN/m] of the release-probability sigmoid. \\
        \bottomrule
    \end{tabular}
\end{table}

\begin{table}[H]
    \caption{Source classification: cortical\_pyramidal\_3comp.}
    \label{tab:prov_cortical_pyramidal_3comp}
    \footnotesize
    \centering
    \begin{tabular}{@{}llp{6.5cm}@{}}
        \toprule
        Parameter & Source & Reference / notes \\
        \midrule
        \texttt{dendrite\_piezo\_scale} & Literature & Lewis \& Grandl 2015; Dalghi 2019 --- Qualitative dendrite-heavy Piezo1 expression \\
        \texttt{g\_Na\_ais} & Literature & Kole 2008 (lower end of 2000-5000 mS/cm\textsuperscript{2} range) --- Default chosen for RK4 stability at dt = 1 $\mu$s; user can raise to 5000 \\
        \texttt{g\_Na\_dend} & Literature & Stuart \& Sakmann 1995 \\
        \texttt{g\_Na\_soma} & Literature & Mainen \& Sejnowski 1996 \\
        \texttt{area\_ratios} & Assumption & Pinsky-Rinzel-style lumped ratios A\_d:A\_s:A\_a = 5:1:0.05 \\
        \texttt{g\_dend\_soma} & Assumption & Lumped axial conductance, mS/cm\textsuperscript{2} of soma \\
        \texttt{g\_soma\_ais} & Assumption & Tighter coupling than dend-soma, reflecting AIS proximity \\
        \bottomrule
    \end{tabular}
\end{table}
 inside supplementary.tex.
% (supplementary.tex preamble already loads \usepackage{float}.)

\begin{table}[H]
    \caption{Source classification: Piezo1.}
    \label{tab:prov_piezo1}
    \footnotesize
    \centering
    \begin{tabular}{@{}llp{6.5cm}@{}}
        \toprule
        Parameter & Source & Reference / notes \\
        \midrule
        \texttt{Ca\_permeability\_fraction} & Literature & Coste et al. 2010 (P\_Ca / P\_Na \textasciitilde{}5) --- \textasciitilde{}10\% of cation current is Ca\textsuperscript{2}\textsuperscript{+} in physiological gradients \\
        \texttt{E\_rev} & Literature & Coste et al. 2010 (non-selective cation, \textasciitilde{}0 mV) \\
        \texttt{T\_half} & Literature & Coste et al. 2010, Science 330:55-60; Lewis \& Grandl 2015 --- Mid-point of activation across multiple cell types \\
        \texttt{delta\_A} & Literature & Coste et al. 2010 (Boltzmann fit to single-channel records) --- Effective gating area change from $\Delta$A$\cdot$T fit \\
        \texttt{g\_max} & Literature & Coste et al. 2010 (single-channel slope) --- Single-channel conductance \textasciitilde{}25 pS in physiological gradients \\
        \texttt{tau\_act} & Literature & Coste et al. 2010 Fig. 2 (activation rise $\tau$ \textasciitilde{}1-2 ms) --- Phase-5 fixed-tau replacement for the bell-shaped $\tau$\_w(T) \\
        \texttt{tau\_inact} & Literature & Coste et al. 2010 Fig. 2 (inactivation $\tau$ 10-15 ms) \\
        \texttt{g\_density} & Assumption & --- --- Sparse expression assumption; user typically overrides per cell type \\
        \texttt{inact\_T\_half} & Assumption & --- --- Set above activation T\_half; not directly fit in Coste 2010 \\
        \texttt{inact\_delta\_A} & Assumption & --- --- Comparable steepness to activation; placeholder pending dedicated fit \\
        \texttt{tau\_0} & Assumption & --- --- Single time constant subsuming activation kinetics; superseded by tau\_act when set \\
        \bottomrule
    \end{tabular}
\end{table}

\begin{table}[H]
    \caption{Source classification: TRAAK.}
    \label{tab:prov_traak}
    \footnotesize
    \centering
    \begin{tabular}{@{}llp{6.5cm}@{}}
        \toprule
        Parameter & Source & Reference / notes \\
        \midrule
        \texttt{E\_rev} & Literature & K\textsuperscript{+} Nernst at physiological [K] --- -90 mV typical \\
        \texttt{T\_half} & Literature & Brohawn et al. 2014, PNAS 111:3614-3619 \\
        \texttt{delta\_A} & Literature & Brohawn et al. 2014 \\
        \texttt{g\_max} & Literature & Brohawn et al. 2014 \\
        \texttt{tau\_act} & Literature & Brohawn et al. 2014 (single-channel mechanical activation rise) --- Phase-5 fixed-tau replacement for the bell-shaped $\tau$\_w(T) \\
        \texttt{g\_density} & Assumption & Estimated from neural expression levels; user-tunable \\
        \texttt{tau\_0} & Assumption & Legacy bell-shape baseline; superseded by tau\_act on the library preset \\
        \texttt{Ca\_permeability\_fraction} & Derived & K2P channels are K\textsuperscript{+}-selective; no Ca\textsuperscript{2}\textsuperscript{+} flux \\
        \texttt{inact\_T\_half} & Derived & Not applicable; tau\_inact is None \\
        \texttt{inact\_delta\_A} & Derived & Not applicable; tau\_inact is None \\
        \texttt{tau\_inact} & Derived & K2P channels don't inactivate; field defaults to None (gate held at 1) \\
        \bottomrule
    \end{tabular}
\end{table}

\begin{table}[H]
    \caption{Source classification: TREK-1.}
    \label{tab:prov_trek_1}
    \footnotesize
    \centering
    \begin{tabular}{@{}llp{6.5cm}@{}}
        \toprule
        Parameter & Source & Reference / notes \\
        \midrule
        \texttt{E\_rev} & Literature & K\textsuperscript{+} Nernst \\
        \texttt{T\_half} & Literature & Honor\'{e} 2007, Nat Rev Neurosci 8:251-261 \\
        \texttt{delta\_A} & Literature & Honor\'{e} 2007 \\
        \texttt{g\_max} & Literature & Honor\'{e} 2007 \\
        \texttt{tau\_act} & Literature & Honor\'{e} 2007 (mid-range of 3-10 ms activation $\tau$ reported in review) --- Phase-5 fixed-tau replacement for the bell-shaped $\tau$\_w(T) \\
        \texttt{g\_density} & Assumption & Cortical pyramidal neuron expression range \\
        \texttt{tau\_0} & Assumption & Legacy bell-shape baseline; superseded by tau\_act on the library preset \\
        \texttt{Ca\_permeability\_fraction} & Derived & K2P channels are K\textsuperscript{+}-selective; no Ca\textsuperscript{2}\textsuperscript{+} flux \\
        \texttt{inact\_T\_half} & Derived & Not applicable; tau\_inact is None \\
        \texttt{inact\_delta\_A} & Derived & Not applicable; tau\_inact is None \\
        \texttt{tau\_inact} & Derived & K2P channels don't inactivate; field defaults to None \\
        \bottomrule
    \end{tabular}
\end{table}

\begin{table}[H]
    \caption{Source classification: TREK-2.}
    \label{tab:prov_trek_2}
    \footnotesize
    \centering
    \begin{tabular}{@{}llp{6.5cm}@{}}
        \toprule
        Parameter & Source & Reference / notes \\
        \midrule
        \texttt{E\_rev} & Literature & K\textsuperscript{+} Nernst \\
        \texttt{T\_half} & Literature & Kang \& Kim 2024 review \\
        \texttt{delta\_A} & Literature & Kang \& Kim 2024 \\
        \texttt{g\_max} & Literature & Kang \& Kim 2024 \\
        \texttt{tau\_act} & Literature & Kang \& Kim 2024 (TREK-2 kinetics; similar to TREK-1) --- Phase-5 fixed-tau replacement for the bell-shaped $\tau$\_w(T) \\
        \texttt{g\_density} & Assumption & Cortical pyramidal neuron expression range \\
        \texttt{tau\_0} & Assumption & Legacy bell-shape baseline; superseded by tau\_act on the library preset \\
        \texttt{Ca\_permeability\_fraction} & Derived & K2P channels are K\textsuperscript{+}-selective; no Ca\textsuperscript{2}\textsuperscript{+} flux \\
        \texttt{inact\_T\_half} & Derived & Not applicable; tau\_inact is None \\
        \texttt{inact\_delta\_A} & Derived & Not applicable; tau\_inact is None \\
        \texttt{tau\_inact} & Derived & K2P channels don't inactivate; field defaults to None \\
        \bottomrule
    \end{tabular}
\end{table}

\begin{table}[H]
    \caption{Source classification: TRPV1.}
    \label{tab:prov_trpv1}
    \footnotesize
    \centering
    \begin{tabular}{@{}llp{6.5cm}@{}}
        \toprule
        Parameter & Source & Reference / notes \\
        \midrule
        \texttt{E\_rev} & Literature & Slightly Ca\textsuperscript{2}\textsuperscript{+}-preferring non-selective cation, \textasciitilde{}10 mV \\
        \texttt{T\_half} & Literature & Caterina et al. 1997, 2000 (heat activation threshold) \\
        \texttt{k} & Calibrated & Boltzmann steepness chosen so Po(47)/Po(37) > 20 matching the literature Q10 $\geq$ 20 floor (Caterina 1997). \\
        \texttt{g\_density} & Assumption & Conservative mid-range density when expressed \\
        \bottomrule
    \end{tabular}
\end{table}

\begin{table}[H]
    \caption{Source classification: TRPV4.}
    \label{tab:prov_trpv4}
    \footnotesize
    \centering
    \begin{tabular}{@{}llp{6.5cm}@{}}
        \toprule
        Parameter & Source & Reference / notes \\
        \midrule
        \texttt{E\_rev} & Literature & Non-selective cation, \textasciitilde{}0 mV \\
        \texttt{T\_half} & Literature & Watanabe et al. 2002 (warm-activation threshold) --- Within the 27-35$^\circ$C literature range \\
        \texttt{k} & Calibrated & Q10 $\approx$ 10 around threshold, less steep than TRPV1 \\
        \texttt{g\_density} & Assumption & Conservative mid-range density when expressed \\
        \bottomrule
    \end{tabular}
\end{table}

\begin{table}[H]
    \caption{Source classification: TRPA1.}
    \label{tab:prov_trpa1}
    \footnotesize
    \centering
    \begin{tabular}{@{}llp{6.5cm}@{}}
        \toprule
        Parameter & Source & Reference / notes \\
        \midrule
        \texttt{Ca\_permeability\_fraction} & Literature & Karashima et al. 2010; Wang et al. 2008 (P\_Ca/P\_Na $\approx$ 5--8) --- \textasciitilde{}15--20\% of TRPA1 cation current is Ca\textsuperscript{2}\textsuperscript{+} at physiological gradients. \\
        \texttt{E\_rev} & Literature & Karashima et al. 2010 (non-selective cation, E\_rev $\approx$ 0 mV) \\
        \texttt{g\_max} & Literature & Karashima et al. 2010; Wang et al. 2008 (TRPA1 single-channel \textasciitilde{}tens of pS) \\
        \texttt{g\_density} & Assumption & Astrocytic TRPA1 surface density; user overrides per preparation. \\
        \texttt{tau\_0} & Assumption & Legacy bell-shape baseline; superseded by tau\_act. \\
        \texttt{T\_half} & Estimate & Oh et al. 2019, Curr Biol 29:3386 (US activates astrocytic TRPA1) --- Half-activation tension for TRPA1 mechanogating is not directly quantified; set between Piezo1 (2.7) and TRAAK (5.0) mN/m as a best guess pending a dedicated tension-clamp fit. \\
        \texttt{delta\_A} & Estimate & Gating-area sensitivity guessed comparable to other MS channels (\textasciitilde{}15 nm\textsuperscript{2}). \\
        \texttt{tau\_act} & Estimate & Channel gating is fast (ms); the slow relay signature lives in the downstream Ca/glutamate dynamics, not in TRPA1 gating. \\
        \bottomrule
    \end{tabular}
\end{table}

\begin{table}[H]
    \caption{Source classification: strain\_to\_tension.}
    \label{tab:prov_strain_to_tension}
    \footnotesize
    \centering
    \begin{tabular}{@{}llp{6.5cm}@{}}
        \toprule
        Parameter & Source & Reference / notes \\
        \midrule
        \texttt{K\_A} & Assumption & Lipid-bilayer area-expansion modulus, range 0.1-0.5 N/m --- Default 0.25 N/m is the centre of the published range (Rawicz 2000, Sukharev 2012); the actual coupling from tissue-scale strain to per-channel membrane tension subsumes geometric concentration, cytoskeletal mechanics, and membrane curvature effects. User must calibrate for the specific tissue/cell-type/exposure under study. \\
        \texttt{coupling\_factor} & Assumption & Dimensionless multi-scale amplification $\alpha$ relating far-field tissue strain to local membrane area change. Code default 1.0; manuscript-canonical $\alpha$=1000 (cortical-pyramidal-calibrated, 5$\times$ the iso-potential/hippocampal literature $\alpha$$\approx$200). Calibrate per tissue/cell-type/exposure. \\
        \bottomrule
    \end{tabular}
\end{table}

\begin{table}[H]
    \caption{Source classification: NICEModel.}
    \label{tab:prov_nicemodel}
    \footnotesize
    \centering
    \begin{tabular}{@{}llp{6.5cm}@{}}
        \toprule
        Parameter & Source & Reference / notes \\
        \midrule
        \texttt{C\_m0} & Literature & Standard neuronal specific capacitance, 1 $\mu$F/cm\textsuperscript{2} \\
        \texttt{delta\_0} & Literature & Lipid bilayer half-thickness (Plaksin 2014) \\
        \texttt{f\_carrier} & Literature & Typical TUS carrier (Legon, Tyler) --- 500 kHz default; user supplies their protocol value \\
        \texttt{f\_resonant} & Literature & Plaksin et al. 2014 (leaflet resonance \textasciitilde{}5 MHz) \\
        \texttt{lj\_strength} & Calibrated & Dimensionless prefactor for the LJ-like wall on Z<0; calibrated so the wall contribution stays <5\% of linear restoring at |Z| < 0.1$\cdot$$\delta$\_0 \\
        \texttt{m\_eff} & Calibrated & SONIC operating range, Lemaire et al. 2019 --- Calibrated so |Z|(100 kPa) $\approx$ 0.6 nm $\to$ $\langle$$\Delta$C\_m$\rangle$/C\_m0 $\approx$ 5\% \\
        \texttt{st\_strength} & Calibrated & Dimensionless prefactor for the surface-tension wall on Z>0; default 1.0 doubles asymptotic restoring on the expansion side, extending the quantitative validity range to \textasciitilde{}1 MPa (Plaksin \& Krasovitski 2014 Fig. 2). \\
        \texttt{damping\_ratio} & Assumption & $\zeta$ = 0.1 (Q $\approx$ 5); Plaksin-consistent moderate damping \\
        \bottomrule
    \end{tabular}
\end{table}

\begin{table}[H]
    \caption{Source classification: SolitonModel.}
    \label{tab:prov_solitonmodel}
    \footnotesize
    \centering
    \begin{tabular}{@{}llp{6.5cm}@{}}
        \toprule
        Parameter & Source & Reference / notes \\
        \midrule
        \texttt{T\_m} & Literature & Heimburg 2007, Thermal Biophysics of Membranes; Heimburg \& Jackson 2005, PNAS 102:9790 --- Membrane melting transition. Biological membranes melt \textasciitilde{}10-15$^\circ$C below body temperature; default 30$^\circ$C is representative (lung-surfactant T\_m $\approx$ 30$^\circ$C, the system Heimburg \& Jackson fit). Must be < 37$^\circ$C for the heating-reduces-excitability sign. \\
        \texttt{c} & Literature & Soft-tissue speed of sound; matches repo FDTD C0 = 1540 m/s --- $\rho$c\textsuperscript{2} is the adiabatic bulk modulus for the strain conversion. \\
        \texttt{rho} & Literature & Soft-tissue/water density; matches repo FDTD RHO0 = 1000 kg/m\textsuperscript{3} --- Used in the bulk-modulus strain $\varepsilon$ = P/($\rho$c\textsuperscript{2}). \\
        \texttt{peak\_shape} & Assumption & Gaussian (default) vs Lorentzian shape for the normalised c\_p peak. \\
        \texttt{c\_p\_max} & Estimate & Heimburg 2007 ($\kappa$\_A\^{}T $\propto$ $\Delta$c\_p proportionality) --- Dimensionless peak-to-baseline excess-compressibility ratio encoding Heimburg's c\_p $\leftrightarrow$ compressibility divergence near T\_m. Tuned amplitude (default 10$\times$). \\
        \texttt{sigma\_T} & Estimate & Heimburg 2007 (broadened multicomponent-membrane transition) --- Half-width of the c\_p / compressibility peak, a few $^\circ$C. Biological transitions are broad; specific value tuned. \\
        \bottomrule
    \end{tabular}
\end{table}

\begin{table}[H]
    \caption{Source classification: CalciumDynamics.}
    \label{tab:prov_calciumdynamics}
    \footnotesize
    \centering
    \begin{tabular}{@{}llp{6.5cm}@{}}
        \toprule
        Parameter & Source & Reference / notes \\
        \midrule
        \texttt{Ca\_rest} & Literature & Cortical pyramidal resting [Ca\textsuperscript{2}\textsuperscript{+}]\_i, \textasciitilde{}50 nM \\
        \texttt{tau\_decay} & Literature & Mainen \& Sejnowski 1996; Borg-Graham 1999 --- Single-pool decay $\tau$ in cortical models, range 30-80 ms \\
        \texttt{flux\_factor} & Calibrated & Lumped sub-membrane shell volume $\times$ buffering fraction $\times$ 1/(z$\cdot$F). Default 0.5 lies within published current-to-concentration coupling factors used in cortical Ca\textsuperscript{2}\textsuperscript{+}-pool models. \\
        \bottomrule
    \end{tabular}
\end{table}

\begin{table}[H]
    \caption{Source classification: KCaChannel.}
    \label{tab:prov_kcachannel}
    \footnotesize
    \centering
    \begin{tabular}{@{}llp{6.5cm}@{}}
        \toprule
        Parameter & Source & Reference / notes \\
        \midrule
        \texttt{Ca\_half} & Literature & K\"{o}hler et al. 1996 (SK channel) \\
        \texttt{E\_K} & Literature & K\textsuperscript{+} Nernst, matches Pospischil RS \\
        \texttt{n\_hill} & Literature & Hirschberg et al. 1998 \\
        \texttt{g\_density} & Assumption & Cortical pyramidal SK density range, 0.1--1.0 mS/cm\textsuperscript{2} (Sah 1996) \\
        \bottomrule
    \end{tabular}
\end{table}

\begin{table}[H]
    \caption{Source classification: AstrocyteRelay.}
    \label{tab:prov_astrocyterelay}
    \footnotesize
    \centering
    \begin{tabular}{@{}llp{6.5cm}@{}}
        \toprule
        Parameter & Source & Reference / notes \\
        \midrule
        \texttt{Ca\_rest} & Literature & Bazargani \& Attwell 2016 (resting astrocytic [Ca\textsuperscript{2}\textsuperscript{+}] \textasciitilde{}50--100 nM) \\
        \texttt{V\_astro} & Literature & Bazargani \& Attwell 2016, Nat Neurosci (astrocyte V\_rest $\approx$ -80 mV) --- Sets the Ca\textsuperscript{2}\textsuperscript{+} driving force through TRPA1 (non-spiking cell). \\
        \texttt{tau\_decay} & Literature & Bazargani \& Attwell 2016; Oh et al. 2019 (astrocytic Ca\textsuperscript{2}\textsuperscript{+} transients decay over 100s of ms) --- Slow relay: hundreds of ms, not per-carrier-cycle. \\
        \texttt{flux\_factor} & Calibrated & Chosen so the astrocytic Ca\textsuperscript{2}\textsuperscript{+} pool peaks \textasciitilde{}1 $\mu$M at full US drive (physiological gliotransmission regime), given TRPA1 g\_density, Ca permeability, and the -80 mV driving force. \\
        \texttt{glu\_Ca\_half} & Estimate & Parpura \& Haydon 2000 (Ca-dependent gliotransmission threshold) --- Half-release [Ca\textsuperscript{2}\textsuperscript{+}] a few hundred nM; tuned for gliotransmission onset. \\
        \texttt{glu\_max} & Estimate & Peak cleft glutamate scale [$\mu$M] from gliotransmission; order-of-magnitude. \\
        \texttt{glu\_n\_hill} & Estimate & Cooperative Ca\textsuperscript{2}\textsuperscript{+}-triggered release; Hill coefficient \textasciitilde{}2--3 (vesicular sensor). \\
        \bottomrule
    \end{tabular}
\end{table}

\begin{table}[H]
    \caption{Source classification: NMDAReceptor.}
    \label{tab:prov_nmdareceptor}
    \footnotesize
    \centering
    \begin{tabular}{@{}llp{6.5cm}@{}}
        \toprule
        Parameter & Source & Reference / notes \\
        \midrule
        \texttt{E\_rev} & Literature & NMDA non-selective cation, E\_rev $\approx$ 0 mV \\
        \texttt{K\_glu} & Literature & Patneau \& Mayer 1990 (NMDA glutamate EC50 \textasciitilde{}1--3 $\mu$M) \\
        \texttt{Mg} & Literature & Jahr \& Stevens 1990 (physiological extracellular [Mg\textsuperscript{2}\textsuperscript{+}] $\approx$ 1 mM) \\
        \texttt{mg\_block\_v\_scale} & Literature & Jahr \& Stevens 1990 (voltage-dependent Mg\textsuperscript{2}\textsuperscript{+} block: 16.13 mV, 3.57 mM) --- g(V) $\propto$ 1/(1 + [Mg]$\cdot$exp(-V/16.13)/3.57). \\
        \texttt{alpha\_NMDA} & Calibrated & Maneshi et al. 2017, Sci Rep (membrane stretch activates NMDARs); Singh et al. 2012 --- Max fractional stretch enhancement of NMDA current. Default 1.0 $\to$ current roughly doubles at saturating stretch, matching Maneshi 2017 (stretch $\approx$ doubled NMDA current). f\_stretch($\infty$) = 1 + alpha\_NMDA. \\
        \texttt{g\_max} & Assumption & Post-synaptic NMDA conductance density; user overrides per cell type. \\
        \texttt{T\_half\_stretch} & Estimate & Maneshi et al. 2017 --- Half-enhancement membrane tension [mN/m]. Mapped to the MS-channel tension range pending a dedicated stretch-clamp NMDA fit. \\
        \texttt{k\_stretch} & Estimate & Boltzmann slope [mN/m] of the stretch sigmoid; comparable steepness to MS channels. \\
        \bottomrule
    \end{tabular}
\end{table}

\begin{table}[H]
    \caption{Source classification: MechanoNav.}
    \label{tab:prov_mechanonav}
    \footnotesize
    \centering
    \begin{tabular}{@{}llp{6.5cm}@{}}
        \toprule
        Parameter & Source & Reference / notes \\
        \midrule
        \texttt{k\_a} & Literature & Hodgkin-Huxley-type Nav activation Boltzmann slope --- Activation slope factor [mV]. \\
        \texttt{v\_half0} & Literature & Representative mammalian Nav activation midpoint (Hodgkin-Huxley-type) --- Baseline activation half-voltage [mV] for the scalar gain. \\
        \texttt{k\_mechano} & Calibrated & Kubanek et al. 2016, Sci Rep (US shifts Nav/Cav voltage dependence) --- Hyperpolarising activation-shift sensitivity [mV per (mN/m)]. K\_MECHANO\_NAV = 0.3 is calibrated so a few-mN/m tension yields a \textasciitilde{}mV-scale activation shift and a tens-of-percent Na\textsuperscript{+} current change, the order Kubanek report (not a directly tabulated value). The neuron field k\_mechano\_nav defaults to 0 (pathway off) to preserve the baseline neuron; set it to K\_MECHANO\_NAV to enable. \\
        \texttt{v\_eval} & Assumption & Near-threshold reference voltage [mV] at which nav\_current\_gain reads the open fraction. \\
        \bottomrule
    \end{tabular}
\end{table}

\begin{table}[H]
    \caption{Source classification: MechanoSynapse.}
    \label{tab:prov_mechanosynapse}
    \footnotesize
    \centering
    \begin{tabular}{@{}llp{6.5cm}@{}}
        \toprule
        Parameter & Source & Reference / notes \\
        \midrule
        \texttt{E\_ampa} & Literature & AMPA non-selective cation reversal, \textasciitilde{}0 mV \\
        \texttt{K\_ampa} & Literature & AMPA is low-affinity (glutamate EC50 \textasciitilde{} hundreds of $\mu$M) --- AMPA glutamate half-saturation [$\mu$M]. \\
        \texttt{Pr\_baseline} & Literature & Tyler 2008; central-synapse release probability range 0.1--0.3 --- Resting vesicle release probability at zero tension. \\
        \texttt{tau\_glu} & Literature & Cleft glutamate clearance / uptake time constant (few ms) --- First-order glutamate pool decay [ms]. \\
        \texttt{g\_ampa} & Calibrated & Post-synaptic AMPA conductance density [mS/cm\textsuperscript{2}]; calibrated with glu\_scale. \\
        \texttt{glu\_scale} & Calibrated & Sustained cleft glutamate [$\mu$M] at full release; calibrated with g\_ampa so the resting (baseline-Pr) state is sub-threshold and US-driven release brings the post neuron to firing. \\
        \texttt{Pr\_max} & Estimate & Tyler 2012 (mechanical perturbation raises release probability) --- Saturating release probability at high presynaptic tension. \\
        \texttt{T\_half\_pre} & Estimate & Tyler 2012 --- Half-modulation presynaptic tension [mN/m]; in the MS-channel range. \\
        \texttt{k\_pre} & Estimate & Boltzmann slope [mN/m] of the release-probability sigmoid. \\
        \bottomrule
    \end{tabular}
\end{table}

\begin{table}[H]
    \caption{Source classification: cortical\_pyramidal\_3comp.}
    \label{tab:prov_cortical_pyramidal_3comp}
    \footnotesize
    \centering
    \begin{tabular}{@{}llp{6.5cm}@{}}
        \toprule
        Parameter & Source & Reference / notes \\
        \midrule
        \texttt{dendrite\_piezo\_scale} & Literature & Lewis \& Grandl 2015; Dalghi 2019 --- Qualitative dendrite-heavy Piezo1 expression \\
        \texttt{g\_Na\_ais} & Literature & Kole 2008 (lower end of 2000-5000 mS/cm\textsuperscript{2} range) --- Default chosen for RK4 stability at dt = 1 $\mu$s; user can raise to 5000 \\
        \texttt{g\_Na\_dend} & Literature & Stuart \& Sakmann 1995 \\
        \texttt{g\_Na\_soma} & Literature & Mainen \& Sejnowski 1996 \\
        \texttt{area\_ratios} & Assumption & Pinsky-Rinzel-style lumped ratios A\_d:A\_s:A\_a = 5:1:0.05 \\
        \texttt{g\_dend\_soma} & Assumption & Lumped axial conductance, mS/cm\textsuperscript{2} of soma \\
        \texttt{g\_soma\_ais} & Assumption & Tighter coupling than dend-soma, reflecting AIS proximity \\
        \bottomrule
    \end{tabular}
\end{table}
 inside supplementary.tex.
% (supplementary.tex preamble already loads \usepackage{float}.)

\begin{table}[H]
    \caption{Source classification: Piezo1.}
    \label{tab:prov_piezo1}
    \footnotesize
    \centering
    \begin{tabular}{@{}llp{6.5cm}@{}}
        \toprule
        Parameter & Source & Reference / notes \\
        \midrule
        \texttt{Ca\_permeability\_fraction} & Literature & Coste et al. 2010 (P\_Ca / P\_Na \textasciitilde{}5) --- \textasciitilde{}10\% of cation current is Ca\textsuperscript{2}\textsuperscript{+} in physiological gradients \\
        \texttt{E\_rev} & Literature & Coste et al. 2010 (non-selective cation, \textasciitilde{}0 mV) \\
        \texttt{T\_half} & Literature & Coste et al. 2010, Science 330:55-60; Lewis \& Grandl 2015 --- Mid-point of activation across multiple cell types \\
        \texttt{delta\_A} & Literature & Coste et al. 2010 (Boltzmann fit to single-channel records) --- Effective gating area change from $\Delta$A$\cdot$T fit \\
        \texttt{g\_max} & Literature & Coste et al. 2010 (single-channel slope) --- Single-channel conductance \textasciitilde{}25 pS in physiological gradients \\
        \texttt{tau\_act} & Literature & Coste et al. 2010 Fig. 2 (activation rise $\tau$ \textasciitilde{}1-2 ms) --- Phase-5 fixed-tau replacement for the bell-shaped $\tau$\_w(T) \\
        \texttt{tau\_inact} & Literature & Coste et al. 2010 Fig. 2 (inactivation $\tau$ 10-15 ms) \\
        \texttt{g\_density} & Assumption & --- --- Sparse expression assumption; user typically overrides per cell type \\
        \texttt{inact\_T\_half} & Assumption & --- --- Set above activation T\_half; not directly fit in Coste 2010 \\
        \texttt{inact\_delta\_A} & Assumption & --- --- Comparable steepness to activation; placeholder pending dedicated fit \\
        \texttt{tau\_0} & Assumption & --- --- Single time constant subsuming activation kinetics; superseded by tau\_act when set \\
        \bottomrule
    \end{tabular}
\end{table}

\begin{table}[H]
    \caption{Source classification: TRAAK.}
    \label{tab:prov_traak}
    \footnotesize
    \centering
    \begin{tabular}{@{}llp{6.5cm}@{}}
        \toprule
        Parameter & Source & Reference / notes \\
        \midrule
        \texttt{E\_rev} & Literature & K\textsuperscript{+} Nernst at physiological [K] --- -90 mV typical \\
        \texttt{T\_half} & Literature & Brohawn et al. 2014, PNAS 111:3614-3619 \\
        \texttt{delta\_A} & Literature & Brohawn et al. 2014 \\
        \texttt{g\_max} & Literature & Brohawn et al. 2014 \\
        \texttt{tau\_act} & Literature & Brohawn et al. 2014 (single-channel mechanical activation rise) --- Phase-5 fixed-tau replacement for the bell-shaped $\tau$\_w(T) \\
        \texttt{g\_density} & Assumption & Estimated from neural expression levels; user-tunable \\
        \texttt{tau\_0} & Assumption & Legacy bell-shape baseline; superseded by tau\_act on the library preset \\
        \texttt{Ca\_permeability\_fraction} & Derived & K2P channels are K\textsuperscript{+}-selective; no Ca\textsuperscript{2}\textsuperscript{+} flux \\
        \texttt{inact\_T\_half} & Derived & Not applicable; tau\_inact is None \\
        \texttt{inact\_delta\_A} & Derived & Not applicable; tau\_inact is None \\
        \texttt{tau\_inact} & Derived & K2P channels don't inactivate; field defaults to None (gate held at 1) \\
        \bottomrule
    \end{tabular}
\end{table}

\begin{table}[H]
    \caption{Source classification: TREK-1.}
    \label{tab:prov_trek_1}
    \footnotesize
    \centering
    \begin{tabular}{@{}llp{6.5cm}@{}}
        \toprule
        Parameter & Source & Reference / notes \\
        \midrule
        \texttt{E\_rev} & Literature & K\textsuperscript{+} Nernst \\
        \texttt{T\_half} & Literature & Honor\'{e} 2007, Nat Rev Neurosci 8:251-261 \\
        \texttt{delta\_A} & Literature & Honor\'{e} 2007 \\
        \texttt{g\_max} & Literature & Honor\'{e} 2007 \\
        \texttt{tau\_act} & Literature & Honor\'{e} 2007 (mid-range of 3-10 ms activation $\tau$ reported in review) --- Phase-5 fixed-tau replacement for the bell-shaped $\tau$\_w(T) \\
        \texttt{g\_density} & Assumption & Cortical pyramidal neuron expression range \\
        \texttt{tau\_0} & Assumption & Legacy bell-shape baseline; superseded by tau\_act on the library preset \\
        \texttt{Ca\_permeability\_fraction} & Derived & K2P channels are K\textsuperscript{+}-selective; no Ca\textsuperscript{2}\textsuperscript{+} flux \\
        \texttt{inact\_T\_half} & Derived & Not applicable; tau\_inact is None \\
        \texttt{inact\_delta\_A} & Derived & Not applicable; tau\_inact is None \\
        \texttt{tau\_inact} & Derived & K2P channels don't inactivate; field defaults to None \\
        \bottomrule
    \end{tabular}
\end{table}

\begin{table}[H]
    \caption{Source classification: TREK-2.}
    \label{tab:prov_trek_2}
    \footnotesize
    \centering
    \begin{tabular}{@{}llp{6.5cm}@{}}
        \toprule
        Parameter & Source & Reference / notes \\
        \midrule
        \texttt{E\_rev} & Literature & K\textsuperscript{+} Nernst \\
        \texttt{T\_half} & Literature & Kang \& Kim 2024 review \\
        \texttt{delta\_A} & Literature & Kang \& Kim 2024 \\
        \texttt{g\_max} & Literature & Kang \& Kim 2024 \\
        \texttt{tau\_act} & Literature & Kang \& Kim 2024 (TREK-2 kinetics; similar to TREK-1) --- Phase-5 fixed-tau replacement for the bell-shaped $\tau$\_w(T) \\
        \texttt{g\_density} & Assumption & Cortical pyramidal neuron expression range \\
        \texttt{tau\_0} & Assumption & Legacy bell-shape baseline; superseded by tau\_act on the library preset \\
        \texttt{Ca\_permeability\_fraction} & Derived & K2P channels are K\textsuperscript{+}-selective; no Ca\textsuperscript{2}\textsuperscript{+} flux \\
        \texttt{inact\_T\_half} & Derived & Not applicable; tau\_inact is None \\
        \texttt{inact\_delta\_A} & Derived & Not applicable; tau\_inact is None \\
        \texttt{tau\_inact} & Derived & K2P channels don't inactivate; field defaults to None \\
        \bottomrule
    \end{tabular}
\end{table}

\begin{table}[H]
    \caption{Source classification: TRPV1.}
    \label{tab:prov_trpv1}
    \footnotesize
    \centering
    \begin{tabular}{@{}llp{6.5cm}@{}}
        \toprule
        Parameter & Source & Reference / notes \\
        \midrule
        \texttt{E\_rev} & Literature & Slightly Ca\textsuperscript{2}\textsuperscript{+}-preferring non-selective cation, \textasciitilde{}10 mV \\
        \texttt{T\_half} & Literature & Caterina et al. 1997, 2000 (heat activation threshold) \\
        \texttt{k} & Calibrated & Boltzmann steepness chosen so Po(47)/Po(37) > 20 matching the literature Q10 $\geq$ 20 floor (Caterina 1997). \\
        \texttt{g\_density} & Assumption & Conservative mid-range density when expressed \\
        \bottomrule
    \end{tabular}
\end{table}

\begin{table}[H]
    \caption{Source classification: TRPV4.}
    \label{tab:prov_trpv4}
    \footnotesize
    \centering
    \begin{tabular}{@{}llp{6.5cm}@{}}
        \toprule
        Parameter & Source & Reference / notes \\
        \midrule
        \texttt{E\_rev} & Literature & Non-selective cation, \textasciitilde{}0 mV \\
        \texttt{T\_half} & Literature & Watanabe et al. 2002 (warm-activation threshold) --- Within the 27-35$^\circ$C literature range \\
        \texttt{k} & Calibrated & Q10 $\approx$ 10 around threshold, less steep than TRPV1 \\
        \texttt{g\_density} & Assumption & Conservative mid-range density when expressed \\
        \bottomrule
    \end{tabular}
\end{table}

\begin{table}[H]
    \caption{Source classification: TRPA1.}
    \label{tab:prov_trpa1}
    \footnotesize
    \centering
    \begin{tabular}{@{}llp{6.5cm}@{}}
        \toprule
        Parameter & Source & Reference / notes \\
        \midrule
        \texttt{Ca\_permeability\_fraction} & Literature & Karashima et al. 2010; Wang et al. 2008 (P\_Ca/P\_Na $\approx$ 5--8) --- \textasciitilde{}15--20\% of TRPA1 cation current is Ca\textsuperscript{2}\textsuperscript{+} at physiological gradients. \\
        \texttt{E\_rev} & Literature & Karashima et al. 2010 (non-selective cation, E\_rev $\approx$ 0 mV) \\
        \texttt{g\_max} & Literature & Karashima et al. 2010; Wang et al. 2008 (TRPA1 single-channel \textasciitilde{}tens of pS) \\
        \texttt{g\_density} & Assumption & Astrocytic TRPA1 surface density; user overrides per preparation. \\
        \texttt{tau\_0} & Assumption & Legacy bell-shape baseline; superseded by tau\_act. \\
        \texttt{T\_half} & Estimate & Oh et al. 2019, Curr Biol 29:3386 (US activates astrocytic TRPA1) --- Half-activation tension for TRPA1 mechanogating is not directly quantified; set between Piezo1 (2.7) and TRAAK (5.0) mN/m as a best guess pending a dedicated tension-clamp fit. \\
        \texttt{delta\_A} & Estimate & Gating-area sensitivity guessed comparable to other MS channels (\textasciitilde{}15 nm\textsuperscript{2}). \\
        \texttt{tau\_act} & Estimate & Channel gating is fast (ms); the slow relay signature lives in the downstream Ca/glutamate dynamics, not in TRPA1 gating. \\
        \bottomrule
    \end{tabular}
\end{table}

\begin{table}[H]
    \caption{Source classification: strain\_to\_tension.}
    \label{tab:prov_strain_to_tension}
    \footnotesize
    \centering
    \begin{tabular}{@{}llp{6.5cm}@{}}
        \toprule
        Parameter & Source & Reference / notes \\
        \midrule
        \texttt{K\_A} & Assumption & Lipid-bilayer area-expansion modulus, range 0.1-0.5 N/m --- Default 0.25 N/m is the centre of the published range (Rawicz 2000, Sukharev 2012); the actual coupling from tissue-scale strain to per-channel membrane tension subsumes geometric concentration, cytoskeletal mechanics, and membrane curvature effects. User must calibrate for the specific tissue/cell-type/exposure under study. \\
        \texttt{coupling\_factor} & Assumption & Dimensionless multi-scale amplification $\alpha$ relating far-field tissue strain to local membrane area change. Code default 1.0; manuscript-canonical $\alpha$=1000 (cortical-pyramidal-calibrated, 5$\times$ the iso-potential/hippocampal literature $\alpha$$\approx$200). Calibrate per tissue/cell-type/exposure. \\
        \bottomrule
    \end{tabular}
\end{table}

\begin{table}[H]
    \caption{Source classification: NICEModel.}
    \label{tab:prov_nicemodel}
    \footnotesize
    \centering
    \begin{tabular}{@{}llp{6.5cm}@{}}
        \toprule
        Parameter & Source & Reference / notes \\
        \midrule
        \texttt{C\_m0} & Literature & Standard neuronal specific capacitance, 1 $\mu$F/cm\textsuperscript{2} \\
        \texttt{delta\_0} & Literature & Lipid bilayer half-thickness (Plaksin 2014) \\
        \texttt{f\_carrier} & Literature & Typical TUS carrier (Legon, Tyler) --- 500 kHz default; user supplies their protocol value \\
        \texttt{f\_resonant} & Literature & Plaksin et al. 2014 (leaflet resonance \textasciitilde{}5 MHz) \\
        \texttt{lj\_strength} & Calibrated & Dimensionless prefactor for the LJ-like wall on Z<0; calibrated so the wall contribution stays <5\% of linear restoring at |Z| < 0.1$\cdot$$\delta$\_0 \\
        \texttt{m\_eff} & Calibrated & SONIC operating range, Lemaire et al. 2019 --- Calibrated so |Z|(100 kPa) $\approx$ 0.6 nm $\to$ $\langle$$\Delta$C\_m$\rangle$/C\_m0 $\approx$ 5\% \\
        \texttt{st\_strength} & Calibrated & Dimensionless prefactor for the surface-tension wall on Z>0; default 1.0 doubles asymptotic restoring on the expansion side, extending the quantitative validity range to \textasciitilde{}1 MPa (Plaksin \& Krasovitski 2014 Fig. 2). \\
        \texttt{damping\_ratio} & Assumption & $\zeta$ = 0.1 (Q $\approx$ 5); Plaksin-consistent moderate damping \\
        \bottomrule
    \end{tabular}
\end{table}

\begin{table}[H]
    \caption{Source classification: SolitonModel.}
    \label{tab:prov_solitonmodel}
    \footnotesize
    \centering
    \begin{tabular}{@{}llp{6.5cm}@{}}
        \toprule
        Parameter & Source & Reference / notes \\
        \midrule
        \texttt{T\_m} & Literature & Heimburg 2007, Thermal Biophysics of Membranes; Heimburg \& Jackson 2005, PNAS 102:9790 --- Membrane melting transition. Biological membranes melt \textasciitilde{}10-15$^\circ$C below body temperature; default 30$^\circ$C is representative (lung-surfactant T\_m $\approx$ 30$^\circ$C, the system Heimburg \& Jackson fit). Must be < 37$^\circ$C for the heating-reduces-excitability sign. \\
        \texttt{c} & Literature & Soft-tissue speed of sound; matches repo FDTD C0 = 1540 m/s --- $\rho$c\textsuperscript{2} is the adiabatic bulk modulus for the strain conversion. \\
        \texttt{rho} & Literature & Soft-tissue/water density; matches repo FDTD RHO0 = 1000 kg/m\textsuperscript{3} --- Used in the bulk-modulus strain $\varepsilon$ = P/($\rho$c\textsuperscript{2}). \\
        \texttt{peak\_shape} & Assumption & Gaussian (default) vs Lorentzian shape for the normalised c\_p peak. \\
        \texttt{c\_p\_max} & Estimate & Heimburg 2007 ($\kappa$\_A\^{}T $\propto$ $\Delta$c\_p proportionality) --- Dimensionless peak-to-baseline excess-compressibility ratio encoding Heimburg's c\_p $\leftrightarrow$ compressibility divergence near T\_m. Tuned amplitude (default 10$\times$). \\
        \texttt{sigma\_T} & Estimate & Heimburg 2007 (broadened multicomponent-membrane transition) --- Half-width of the c\_p / compressibility peak, a few $^\circ$C. Biological transitions are broad; specific value tuned. \\
        \bottomrule
    \end{tabular}
\end{table}

\begin{table}[H]
    \caption{Source classification: CalciumDynamics.}
    \label{tab:prov_calciumdynamics}
    \footnotesize
    \centering
    \begin{tabular}{@{}llp{6.5cm}@{}}
        \toprule
        Parameter & Source & Reference / notes \\
        \midrule
        \texttt{Ca\_rest} & Literature & Cortical pyramidal resting [Ca\textsuperscript{2}\textsuperscript{+}]\_i, \textasciitilde{}50 nM \\
        \texttt{tau\_decay} & Literature & Mainen \& Sejnowski 1996; Borg-Graham 1999 --- Single-pool decay $\tau$ in cortical models, range 30-80 ms \\
        \texttt{flux\_factor} & Calibrated & Lumped sub-membrane shell volume $\times$ buffering fraction $\times$ 1/(z$\cdot$F). Default 0.5 lies within published current-to-concentration coupling factors used in cortical Ca\textsuperscript{2}\textsuperscript{+}-pool models. \\
        \bottomrule
    \end{tabular}
\end{table}

\begin{table}[H]
    \caption{Source classification: KCaChannel.}
    \label{tab:prov_kcachannel}
    \footnotesize
    \centering
    \begin{tabular}{@{}llp{6.5cm}@{}}
        \toprule
        Parameter & Source & Reference / notes \\
        \midrule
        \texttt{Ca\_half} & Literature & K\"{o}hler et al. 1996 (SK channel) \\
        \texttt{E\_K} & Literature & K\textsuperscript{+} Nernst, matches Pospischil RS \\
        \texttt{n\_hill} & Literature & Hirschberg et al. 1998 \\
        \texttt{g\_density} & Assumption & Cortical pyramidal SK density range, 0.1--1.0 mS/cm\textsuperscript{2} (Sah 1996) \\
        \bottomrule
    \end{tabular}
\end{table}

\begin{table}[H]
    \caption{Source classification: AstrocyteRelay.}
    \label{tab:prov_astrocyterelay}
    \footnotesize
    \centering
    \begin{tabular}{@{}llp{6.5cm}@{}}
        \toprule
        Parameter & Source & Reference / notes \\
        \midrule
        \texttt{Ca\_rest} & Literature & Bazargani \& Attwell 2016 (resting astrocytic [Ca\textsuperscript{2}\textsuperscript{+}] \textasciitilde{}50--100 nM) \\
        \texttt{V\_astro} & Literature & Bazargani \& Attwell 2016, Nat Neurosci (astrocyte V\_rest $\approx$ -80 mV) --- Sets the Ca\textsuperscript{2}\textsuperscript{+} driving force through TRPA1 (non-spiking cell). \\
        \texttt{tau\_decay} & Literature & Bazargani \& Attwell 2016; Oh et al. 2019 (astrocytic Ca\textsuperscript{2}\textsuperscript{+} transients decay over 100s of ms) --- Slow relay: hundreds of ms, not per-carrier-cycle. \\
        \texttt{flux\_factor} & Calibrated & Chosen so the astrocytic Ca\textsuperscript{2}\textsuperscript{+} pool peaks \textasciitilde{}1 $\mu$M at full US drive (physiological gliotransmission regime), given TRPA1 g\_density, Ca permeability, and the -80 mV driving force. \\
        \texttt{glu\_Ca\_half} & Estimate & Parpura \& Haydon 2000 (Ca-dependent gliotransmission threshold) --- Half-release [Ca\textsuperscript{2}\textsuperscript{+}] a few hundred nM; tuned for gliotransmission onset. \\
        \texttt{glu\_max} & Estimate & Peak cleft glutamate scale [$\mu$M] from gliotransmission; order-of-magnitude. \\
        \texttt{glu\_n\_hill} & Estimate & Cooperative Ca\textsuperscript{2}\textsuperscript{+}-triggered release; Hill coefficient \textasciitilde{}2--3 (vesicular sensor). \\
        \bottomrule
    \end{tabular}
\end{table}

\begin{table}[H]
    \caption{Source classification: NMDAReceptor.}
    \label{tab:prov_nmdareceptor}
    \footnotesize
    \centering
    \begin{tabular}{@{}llp{6.5cm}@{}}
        \toprule
        Parameter & Source & Reference / notes \\
        \midrule
        \texttt{E\_rev} & Literature & NMDA non-selective cation, E\_rev $\approx$ 0 mV \\
        \texttt{K\_glu} & Literature & Patneau \& Mayer 1990 (NMDA glutamate EC50 \textasciitilde{}1--3 $\mu$M) \\
        \texttt{Mg} & Literature & Jahr \& Stevens 1990 (physiological extracellular [Mg\textsuperscript{2}\textsuperscript{+}] $\approx$ 1 mM) \\
        \texttt{mg\_block\_v\_scale} & Literature & Jahr \& Stevens 1990 (voltage-dependent Mg\textsuperscript{2}\textsuperscript{+} block: 16.13 mV, 3.57 mM) --- g(V) $\propto$ 1/(1 + [Mg]$\cdot$exp(-V/16.13)/3.57). \\
        \texttt{alpha\_NMDA} & Calibrated & Maneshi et al. 2017, Sci Rep (membrane stretch activates NMDARs); Singh et al. 2012 --- Max fractional stretch enhancement of NMDA current. Default 1.0 $\to$ current roughly doubles at saturating stretch, matching Maneshi 2017 (stretch $\approx$ doubled NMDA current). f\_stretch($\infty$) = 1 + alpha\_NMDA. \\
        \texttt{g\_max} & Assumption & Post-synaptic NMDA conductance density; user overrides per cell type. \\
        \texttt{T\_half\_stretch} & Estimate & Maneshi et al. 2017 --- Half-enhancement membrane tension [mN/m]. Mapped to the MS-channel tension range pending a dedicated stretch-clamp NMDA fit. \\
        \texttt{k\_stretch} & Estimate & Boltzmann slope [mN/m] of the stretch sigmoid; comparable steepness to MS channels. \\
        \bottomrule
    \end{tabular}
\end{table}

\begin{table}[H]
    \caption{Source classification: MechanoNav.}
    \label{tab:prov_mechanonav}
    \footnotesize
    \centering
    \begin{tabular}{@{}llp{6.5cm}@{}}
        \toprule
        Parameter & Source & Reference / notes \\
        \midrule
        \texttt{k\_a} & Literature & Hodgkin-Huxley-type Nav activation Boltzmann slope --- Activation slope factor [mV]. \\
        \texttt{v\_half0} & Literature & Representative mammalian Nav activation midpoint (Hodgkin-Huxley-type) --- Baseline activation half-voltage [mV] for the scalar gain. \\
        \texttt{k\_mechano} & Calibrated & Kubanek et al. 2016, Sci Rep (US shifts Nav/Cav voltage dependence) --- Hyperpolarising activation-shift sensitivity [mV per (mN/m)]. K\_MECHANO\_NAV = 0.3 is calibrated so a few-mN/m tension yields a \textasciitilde{}mV-scale activation shift and a tens-of-percent Na\textsuperscript{+} current change, the order Kubanek report (not a directly tabulated value). The neuron field k\_mechano\_nav defaults to 0 (pathway off) to preserve the baseline neuron; set it to K\_MECHANO\_NAV to enable. \\
        \texttt{v\_eval} & Assumption & Near-threshold reference voltage [mV] at which nav\_current\_gain reads the open fraction. \\
        \bottomrule
    \end{tabular}
\end{table}

\begin{table}[H]
    \caption{Source classification: MechanoSynapse.}
    \label{tab:prov_mechanosynapse}
    \footnotesize
    \centering
    \begin{tabular}{@{}llp{6.5cm}@{}}
        \toprule
        Parameter & Source & Reference / notes \\
        \midrule
        \texttt{E\_ampa} & Literature & AMPA non-selective cation reversal, \textasciitilde{}0 mV \\
        \texttt{K\_ampa} & Literature & AMPA is low-affinity (glutamate EC50 \textasciitilde{} hundreds of $\mu$M) --- AMPA glutamate half-saturation [$\mu$M]. \\
        \texttt{Pr\_baseline} & Literature & Tyler 2008; central-synapse release probability range 0.1--0.3 --- Resting vesicle release probability at zero tension. \\
        \texttt{tau\_glu} & Literature & Cleft glutamate clearance / uptake time constant (few ms) --- First-order glutamate pool decay [ms]. \\
        \texttt{g\_ampa} & Calibrated & Post-synaptic AMPA conductance density [mS/cm\textsuperscript{2}]; calibrated with glu\_scale. \\
        \texttt{glu\_scale} & Calibrated & Sustained cleft glutamate [$\mu$M] at full release; calibrated with g\_ampa so the resting (baseline-Pr) state is sub-threshold and US-driven release brings the post neuron to firing. \\
        \texttt{Pr\_max} & Estimate & Tyler 2012 (mechanical perturbation raises release probability) --- Saturating release probability at high presynaptic tension. \\
        \texttt{T\_half\_pre} & Estimate & Tyler 2012 --- Half-modulation presynaptic tension [mN/m]; in the MS-channel range. \\
        \texttt{k\_pre} & Estimate & Boltzmann slope [mN/m] of the release-probability sigmoid. \\
        \bottomrule
    \end{tabular}
\end{table}

\begin{table}[H]
    \caption{Source classification: cortical\_pyramidal\_3comp.}
    \label{tab:prov_cortical_pyramidal_3comp}
    \footnotesize
    \centering
    \begin{tabular}{@{}llp{6.5cm}@{}}
        \toprule
        Parameter & Source & Reference / notes \\
        \midrule
        \texttt{dendrite\_piezo\_scale} & Literature & Lewis \& Grandl 2015; Dalghi 2019 --- Qualitative dendrite-heavy Piezo1 expression \\
        \texttt{g\_Na\_ais} & Literature & Kole 2008 (lower end of 2000-5000 mS/cm\textsuperscript{2} range) --- Default chosen for RK4 stability at dt = 1 $\mu$s; user can raise to 5000 \\
        \texttt{g\_Na\_dend} & Literature & Stuart \& Sakmann 1995 \\
        \texttt{g\_Na\_soma} & Literature & Mainen \& Sejnowski 1996 \\
        \texttt{area\_ratios} & Assumption & Pinsky-Rinzel-style lumped ratios A\_d:A\_s:A\_a = 5:1:0.05 \\
        \texttt{g\_dend\_soma} & Assumption & Lumped axial conductance, mS/cm\textsuperscript{2} of soma \\
        \texttt{g\_soma\_ais} & Assumption & Tighter coupling than dend-soma, reflecting AIS proximity \\
        \bottomrule
    \end{tabular}
\end{table}
 inside supplementary.tex.
% (supplementary.tex preamble already loads \usepackage{float}.)

\begin{table}[H]
    \caption{Source classification: Piezo1.}
    \label{tab:prov_piezo1}
    \footnotesize
    \centering
    \begin{tabular}{@{}llp{6.5cm}@{}}
        \toprule
        Parameter & Source & Reference / notes \\
        \midrule
        \texttt{Ca\_permeability\_fraction} & Literature & Coste et al. 2010 (P\_Ca / P\_Na \textasciitilde{}5) --- \textasciitilde{}10\% of cation current is Ca\textsuperscript{2}\textsuperscript{+} in physiological gradients \\
        \texttt{E\_rev} & Literature & Coste et al. 2010 (non-selective cation, \textasciitilde{}0 mV) \\
        \texttt{T\_half} & Literature & Coste et al. 2010, Science 330:55-60; Lewis \& Grandl 2015 --- Mid-point of activation across multiple cell types \\
        \texttt{delta\_A} & Literature & Coste et al. 2010 (Boltzmann fit to single-channel records) --- Effective gating area change from $\Delta$A$\cdot$T fit \\
        \texttt{g\_max} & Literature & Coste et al. 2010 (single-channel slope) --- Single-channel conductance \textasciitilde{}25 pS in physiological gradients \\
        \texttt{tau\_act} & Literature & Coste et al. 2010 Fig. 2 (activation rise $\tau$ \textasciitilde{}1-2 ms) --- Phase-5 fixed-tau replacement for the bell-shaped $\tau$\_w(T) \\
        \texttt{tau\_inact} & Literature & Coste et al. 2010 Fig. 2 (inactivation $\tau$ 10-15 ms) \\
        \texttt{g\_density} & Assumption & --- --- Sparse expression assumption; user typically overrides per cell type \\
        \texttt{inact\_T\_half} & Assumption & --- --- Set above activation T\_half; not directly fit in Coste 2010 \\
        \texttt{inact\_delta\_A} & Assumption & --- --- Comparable steepness to activation; placeholder pending dedicated fit \\
        \texttt{tau\_0} & Assumption & --- --- Single time constant subsuming activation kinetics; superseded by tau\_act when set \\
        \bottomrule
    \end{tabular}
\end{table}

\begin{table}[H]
    \caption{Source classification: TRAAK.}
    \label{tab:prov_traak}
    \footnotesize
    \centering
    \begin{tabular}{@{}llp{6.5cm}@{}}
        \toprule
        Parameter & Source & Reference / notes \\
        \midrule
        \texttt{E\_rev} & Literature & K\textsuperscript{+} Nernst at physiological [K] --- -90 mV typical \\
        \texttt{T\_half} & Literature & Brohawn et al. 2014, PNAS 111:3614-3619 \\
        \texttt{delta\_A} & Literature & Brohawn et al. 2014 \\
        \texttt{g\_max} & Literature & Brohawn et al. 2014 \\
        \texttt{tau\_act} & Literature & Brohawn et al. 2014 (single-channel mechanical activation rise) --- Phase-5 fixed-tau replacement for the bell-shaped $\tau$\_w(T) \\
        \texttt{g\_density} & Assumption & Estimated from neural expression levels; user-tunable \\
        \texttt{tau\_0} & Assumption & Legacy bell-shape baseline; superseded by tau\_act on the library preset \\
        \texttt{Ca\_permeability\_fraction} & Derived & K2P channels are K\textsuperscript{+}-selective; no Ca\textsuperscript{2}\textsuperscript{+} flux \\
        \texttt{inact\_T\_half} & Derived & Not applicable; tau\_inact is None \\
        \texttt{inact\_delta\_A} & Derived & Not applicable; tau\_inact is None \\
        \texttt{tau\_inact} & Derived & K2P channels don't inactivate; field defaults to None (gate held at 1) \\
        \bottomrule
    \end{tabular}
\end{table}

\begin{table}[H]
    \caption{Source classification: TREK-1.}
    \label{tab:prov_trek_1}
    \footnotesize
    \centering
    \begin{tabular}{@{}llp{6.5cm}@{}}
        \toprule
        Parameter & Source & Reference / notes \\
        \midrule
        \texttt{E\_rev} & Literature & K\textsuperscript{+} Nernst \\
        \texttt{T\_half} & Literature & Honor\'{e} 2007, Nat Rev Neurosci 8:251-261 \\
        \texttt{delta\_A} & Literature & Honor\'{e} 2007 \\
        \texttt{g\_max} & Literature & Honor\'{e} 2007 \\
        \texttt{tau\_act} & Literature & Honor\'{e} 2007 (mid-range of 3-10 ms activation $\tau$ reported in review) --- Phase-5 fixed-tau replacement for the bell-shaped $\tau$\_w(T) \\
        \texttt{g\_density} & Assumption & Cortical pyramidal neuron expression range \\
        \texttt{tau\_0} & Assumption & Legacy bell-shape baseline; superseded by tau\_act on the library preset \\
        \texttt{Ca\_permeability\_fraction} & Derived & K2P channels are K\textsuperscript{+}-selective; no Ca\textsuperscript{2}\textsuperscript{+} flux \\
        \texttt{inact\_T\_half} & Derived & Not applicable; tau\_inact is None \\
        \texttt{inact\_delta\_A} & Derived & Not applicable; tau\_inact is None \\
        \texttt{tau\_inact} & Derived & K2P channels don't inactivate; field defaults to None \\
        \bottomrule
    \end{tabular}
\end{table}

\begin{table}[H]
    \caption{Source classification: TREK-2.}
    \label{tab:prov_trek_2}
    \footnotesize
    \centering
    \begin{tabular}{@{}llp{6.5cm}@{}}
        \toprule
        Parameter & Source & Reference / notes \\
        \midrule
        \texttt{E\_rev} & Literature & K\textsuperscript{+} Nernst \\
        \texttt{T\_half} & Literature & Kang \& Kim 2024 review \\
        \texttt{delta\_A} & Literature & Kang \& Kim 2024 \\
        \texttt{g\_max} & Literature & Kang \& Kim 2024 \\
        \texttt{tau\_act} & Literature & Kang \& Kim 2024 (TREK-2 kinetics; similar to TREK-1) --- Phase-5 fixed-tau replacement for the bell-shaped $\tau$\_w(T) \\
        \texttt{g\_density} & Assumption & Cortical pyramidal neuron expression range \\
        \texttt{tau\_0} & Assumption & Legacy bell-shape baseline; superseded by tau\_act on the library preset \\
        \texttt{Ca\_permeability\_fraction} & Derived & K2P channels are K\textsuperscript{+}-selective; no Ca\textsuperscript{2}\textsuperscript{+} flux \\
        \texttt{inact\_T\_half} & Derived & Not applicable; tau\_inact is None \\
        \texttt{inact\_delta\_A} & Derived & Not applicable; tau\_inact is None \\
        \texttt{tau\_inact} & Derived & K2P channels don't inactivate; field defaults to None \\
        \bottomrule
    \end{tabular}
\end{table}

\begin{table}[H]
    \caption{Source classification: TRPV1.}
    \label{tab:prov_trpv1}
    \footnotesize
    \centering
    \begin{tabular}{@{}llp{6.5cm}@{}}
        \toprule
        Parameter & Source & Reference / notes \\
        \midrule
        \texttt{E\_rev} & Literature & Slightly Ca\textsuperscript{2}\textsuperscript{+}-preferring non-selective cation, \textasciitilde{}10 mV \\
        \texttt{T\_half} & Literature & Caterina et al. 1997, 2000 (heat activation threshold) \\
        \texttt{k} & Calibrated & Boltzmann steepness chosen so Po(47)/Po(37) > 20 matching the literature Q10 $\geq$ 20 floor (Caterina 1997). \\
        \texttt{g\_density} & Assumption & Conservative mid-range density when expressed \\
        \bottomrule
    \end{tabular}
\end{table}

\begin{table}[H]
    \caption{Source classification: TRPV4.}
    \label{tab:prov_trpv4}
    \footnotesize
    \centering
    \begin{tabular}{@{}llp{6.5cm}@{}}
        \toprule
        Parameter & Source & Reference / notes \\
        \midrule
        \texttt{E\_rev} & Literature & Non-selective cation, \textasciitilde{}0 mV \\
        \texttt{T\_half} & Literature & Watanabe et al. 2002 (warm-activation threshold) --- Within the 27-35$^\circ$C literature range \\
        \texttt{k} & Calibrated & Q10 $\approx$ 10 around threshold, less steep than TRPV1 \\
        \texttt{g\_density} & Assumption & Conservative mid-range density when expressed \\
        \bottomrule
    \end{tabular}
\end{table}

\begin{table}[H]
    \caption{Source classification: TRPA1.}
    \label{tab:prov_trpa1}
    \footnotesize
    \centering
    \begin{tabular}{@{}llp{6.5cm}@{}}
        \toprule
        Parameter & Source & Reference / notes \\
        \midrule
        \texttt{Ca\_permeability\_fraction} & Literature & Karashima et al. 2010; Wang et al. 2008 (P\_Ca/P\_Na $\approx$ 5--8) --- \textasciitilde{}15--20\% of TRPA1 cation current is Ca\textsuperscript{2}\textsuperscript{+} at physiological gradients. \\
        \texttt{E\_rev} & Literature & Karashima et al. 2010 (non-selective cation, E\_rev $\approx$ 0 mV) \\
        \texttt{g\_max} & Literature & Karashima et al. 2010; Wang et al. 2008 (TRPA1 single-channel \textasciitilde{}tens of pS) \\
        \texttt{g\_density} & Assumption & Astrocytic TRPA1 surface density; user overrides per preparation. \\
        \texttt{tau\_0} & Assumption & Legacy bell-shape baseline; superseded by tau\_act. \\
        \texttt{T\_half} & Estimate & Oh et al. 2019, Curr Biol 29:3386 (US activates astrocytic TRPA1) --- Half-activation tension for TRPA1 mechanogating is not directly quantified; set between Piezo1 (2.7) and TRAAK (5.0) mN/m as a best guess pending a dedicated tension-clamp fit. \\
        \texttt{delta\_A} & Estimate & Gating-area sensitivity guessed comparable to other MS channels (\textasciitilde{}15 nm\textsuperscript{2}). \\
        \texttt{tau\_act} & Estimate & Channel gating is fast (ms); the slow relay signature lives in the downstream Ca/glutamate dynamics, not in TRPA1 gating. \\
        \bottomrule
    \end{tabular}
\end{table}

\begin{table}[H]
    \caption{Source classification: strain\_to\_tension.}
    \label{tab:prov_strain_to_tension}
    \footnotesize
    \centering
    \begin{tabular}{@{}llp{6.5cm}@{}}
        \toprule
        Parameter & Source & Reference / notes \\
        \midrule
        \texttt{K\_A} & Assumption & Lipid-bilayer area-expansion modulus, range 0.1-0.5 N/m --- Default 0.25 N/m is the centre of the published range (Rawicz 2000, Sukharev 2012); the actual coupling from tissue-scale strain to per-channel membrane tension subsumes geometric concentration, cytoskeletal mechanics, and membrane curvature effects. User must calibrate for the specific tissue/cell-type/exposure under study. \\
        \texttt{coupling\_factor} & Assumption & Dimensionless multi-scale amplification $\alpha$ relating far-field tissue strain to local membrane area change. Code default 1.0; manuscript-canonical $\alpha$=1000 (cortical-pyramidal-calibrated, 5$\times$ the iso-potential/hippocampal literature $\alpha$$\approx$200). Calibrate per tissue/cell-type/exposure. \\
        \bottomrule
    \end{tabular}
\end{table}

\begin{table}[H]
    \caption{Source classification: NICEModel.}
    \label{tab:prov_nicemodel}
    \footnotesize
    \centering
    \begin{tabular}{@{}llp{6.5cm}@{}}
        \toprule
        Parameter & Source & Reference / notes \\
        \midrule
        \texttt{C\_m0} & Literature & Standard neuronal specific capacitance, 1 $\mu$F/cm\textsuperscript{2} \\
        \texttt{delta\_0} & Literature & Lipid bilayer half-thickness (Plaksin 2014) \\
        \texttt{f\_carrier} & Literature & Typical TUS carrier (Legon, Tyler) --- 500 kHz default; user supplies their protocol value \\
        \texttt{f\_resonant} & Literature & Plaksin et al. 2014 (leaflet resonance \textasciitilde{}5 MHz) \\
        \texttt{lj\_strength} & Calibrated & Dimensionless prefactor for the LJ-like wall on Z<0; calibrated so the wall contribution stays <5\% of linear restoring at |Z| < 0.1$\cdot$$\delta$\_0 \\
        \texttt{m\_eff} & Calibrated & SONIC operating range, Lemaire et al. 2019 --- Calibrated so |Z|(100 kPa) $\approx$ 0.6 nm $\to$ $\langle$$\Delta$C\_m$\rangle$/C\_m0 $\approx$ 5\% \\
        \texttt{st\_strength} & Calibrated & Dimensionless prefactor for the surface-tension wall on Z>0; default 1.0 doubles asymptotic restoring on the expansion side, extending the quantitative validity range to \textasciitilde{}1 MPa (Plaksin \& Krasovitski 2014 Fig. 2). \\
        \texttt{damping\_ratio} & Assumption & $\zeta$ = 0.1 (Q $\approx$ 5); Plaksin-consistent moderate damping \\
        \bottomrule
    \end{tabular}
\end{table}

\begin{table}[H]
    \caption{Source classification: SolitonModel.}
    \label{tab:prov_solitonmodel}
    \footnotesize
    \centering
    \begin{tabular}{@{}llp{6.5cm}@{}}
        \toprule
        Parameter & Source & Reference / notes \\
        \midrule
        \texttt{T\_m} & Literature & Heimburg 2007, Thermal Biophysics of Membranes; Heimburg \& Jackson 2005, PNAS 102:9790 --- Membrane melting transition. Biological membranes melt \textasciitilde{}10-15$^\circ$C below body temperature; default 30$^\circ$C is representative (lung-surfactant T\_m $\approx$ 30$^\circ$C, the system Heimburg \& Jackson fit). Must be < 37$^\circ$C for the heating-reduces-excitability sign. \\
        \texttt{c} & Literature & Soft-tissue speed of sound; matches repo FDTD C0 = 1540 m/s --- $\rho$c\textsuperscript{2} is the adiabatic bulk modulus for the strain conversion. \\
        \texttt{rho} & Literature & Soft-tissue/water density; matches repo FDTD RHO0 = 1000 kg/m\textsuperscript{3} --- Used in the bulk-modulus strain $\varepsilon$ = P/($\rho$c\textsuperscript{2}). \\
        \texttt{peak\_shape} & Assumption & Gaussian (default) vs Lorentzian shape for the normalised c\_p peak. \\
        \texttt{c\_p\_max} & Estimate & Heimburg 2007 ($\kappa$\_A\^{}T $\propto$ $\Delta$c\_p proportionality) --- Dimensionless peak-to-baseline excess-compressibility ratio encoding Heimburg's c\_p $\leftrightarrow$ compressibility divergence near T\_m. Tuned amplitude (default 10$\times$). \\
        \texttt{sigma\_T} & Estimate & Heimburg 2007 (broadened multicomponent-membrane transition) --- Half-width of the c\_p / compressibility peak, a few $^\circ$C. Biological transitions are broad; specific value tuned. \\
        \bottomrule
    \end{tabular}
\end{table}

\begin{table}[H]
    \caption{Source classification: CalciumDynamics.}
    \label{tab:prov_calciumdynamics}
    \footnotesize
    \centering
    \begin{tabular}{@{}llp{6.5cm}@{}}
        \toprule
        Parameter & Source & Reference / notes \\
        \midrule
        \texttt{Ca\_rest} & Literature & Cortical pyramidal resting [Ca\textsuperscript{2}\textsuperscript{+}]\_i, \textasciitilde{}50 nM \\
        \texttt{tau\_decay} & Literature & Mainen \& Sejnowski 1996; Borg-Graham 1999 --- Single-pool decay $\tau$ in cortical models, range 30-80 ms \\
        \texttt{flux\_factor} & Calibrated & Lumped sub-membrane shell volume $\times$ buffering fraction $\times$ 1/(z$\cdot$F). Default 0.5 lies within published current-to-concentration coupling factors used in cortical Ca\textsuperscript{2}\textsuperscript{+}-pool models. \\
        \bottomrule
    \end{tabular}
\end{table}

\begin{table}[H]
    \caption{Source classification: KCaChannel.}
    \label{tab:prov_kcachannel}
    \footnotesize
    \centering
    \begin{tabular}{@{}llp{6.5cm}@{}}
        \toprule
        Parameter & Source & Reference / notes \\
        \midrule
        \texttt{Ca\_half} & Literature & K\"{o}hler et al. 1996 (SK channel) \\
        \texttt{E\_K} & Literature & K\textsuperscript{+} Nernst, matches Pospischil RS \\
        \texttt{n\_hill} & Literature & Hirschberg et al. 1998 \\
        \texttt{g\_density} & Assumption & Cortical pyramidal SK density range, 0.1--1.0 mS/cm\textsuperscript{2} (Sah 1996) \\
        \bottomrule
    \end{tabular}
\end{table}

\begin{table}[H]
    \caption{Source classification: AstrocyteRelay.}
    \label{tab:prov_astrocyterelay}
    \footnotesize
    \centering
    \begin{tabular}{@{}llp{6.5cm}@{}}
        \toprule
        Parameter & Source & Reference / notes \\
        \midrule
        \texttt{Ca\_rest} & Literature & Bazargani \& Attwell 2016 (resting astrocytic [Ca\textsuperscript{2}\textsuperscript{+}] \textasciitilde{}50--100 nM) \\
        \texttt{V\_astro} & Literature & Bazargani \& Attwell 2016, Nat Neurosci (astrocyte V\_rest $\approx$ -80 mV) --- Sets the Ca\textsuperscript{2}\textsuperscript{+} driving force through TRPA1 (non-spiking cell). \\
        \texttt{tau\_decay} & Literature & Bazargani \& Attwell 2016; Oh et al. 2019 (astrocytic Ca\textsuperscript{2}\textsuperscript{+} transients decay over 100s of ms) --- Slow relay: hundreds of ms, not per-carrier-cycle. \\
        \texttt{flux\_factor} & Calibrated & Chosen so the astrocytic Ca\textsuperscript{2}\textsuperscript{+} pool peaks \textasciitilde{}1 $\mu$M at full US drive (physiological gliotransmission regime), given TRPA1 g\_density, Ca permeability, and the -80 mV driving force. \\
        \texttt{glu\_Ca\_half} & Estimate & Parpura \& Haydon 2000 (Ca-dependent gliotransmission threshold) --- Half-release [Ca\textsuperscript{2}\textsuperscript{+}] a few hundred nM; tuned for gliotransmission onset. \\
        \texttt{glu\_max} & Estimate & Peak cleft glutamate scale [$\mu$M] from gliotransmission; order-of-magnitude. \\
        \texttt{glu\_n\_hill} & Estimate & Cooperative Ca\textsuperscript{2}\textsuperscript{+}-triggered release; Hill coefficient \textasciitilde{}2--3 (vesicular sensor). \\
        \bottomrule
    \end{tabular}
\end{table}

\begin{table}[H]
    \caption{Source classification: NMDAReceptor.}
    \label{tab:prov_nmdareceptor}
    \footnotesize
    \centering
    \begin{tabular}{@{}llp{6.5cm}@{}}
        \toprule
        Parameter & Source & Reference / notes \\
        \midrule
        \texttt{E\_rev} & Literature & NMDA non-selective cation, E\_rev $\approx$ 0 mV \\
        \texttt{K\_glu} & Literature & Patneau \& Mayer 1990 (NMDA glutamate EC50 \textasciitilde{}1--3 $\mu$M) \\
        \texttt{Mg} & Literature & Jahr \& Stevens 1990 (physiological extracellular [Mg\textsuperscript{2}\textsuperscript{+}] $\approx$ 1 mM) \\
        \texttt{mg\_block\_v\_scale} & Literature & Jahr \& Stevens 1990 (voltage-dependent Mg\textsuperscript{2}\textsuperscript{+} block: 16.13 mV, 3.57 mM) --- g(V) $\propto$ 1/(1 + [Mg]$\cdot$exp(-V/16.13)/3.57). \\
        \texttt{alpha\_NMDA} & Calibrated & Maneshi et al. 2017, Sci Rep (membrane stretch activates NMDARs); Singh et al. 2012 --- Max fractional stretch enhancement of NMDA current. Default 1.0 $\to$ current roughly doubles at saturating stretch, matching Maneshi 2017 (stretch $\approx$ doubled NMDA current). f\_stretch($\infty$) = 1 + alpha\_NMDA. \\
        \texttt{g\_max} & Assumption & Post-synaptic NMDA conductance density; user overrides per cell type. \\
        \texttt{T\_half\_stretch} & Estimate & Maneshi et al. 2017 --- Half-enhancement membrane tension [mN/m]. Mapped to the MS-channel tension range pending a dedicated stretch-clamp NMDA fit. \\
        \texttt{k\_stretch} & Estimate & Boltzmann slope [mN/m] of the stretch sigmoid; comparable steepness to MS channels. \\
        \bottomrule
    \end{tabular}
\end{table}

\begin{table}[H]
    \caption{Source classification: MechanoNav.}
    \label{tab:prov_mechanonav}
    \footnotesize
    \centering
    \begin{tabular}{@{}llp{6.5cm}@{}}
        \toprule
        Parameter & Source & Reference / notes \\
        \midrule
        \texttt{k\_a} & Literature & Hodgkin-Huxley-type Nav activation Boltzmann slope --- Activation slope factor [mV]. \\
        \texttt{v\_half0} & Literature & Representative mammalian Nav activation midpoint (Hodgkin-Huxley-type) --- Baseline activation half-voltage [mV] for the scalar gain. \\
        \texttt{k\_mechano} & Calibrated & Kubanek et al. 2016, Sci Rep (US shifts Nav/Cav voltage dependence) --- Hyperpolarising activation-shift sensitivity [mV per (mN/m)]. K\_MECHANO\_NAV = 0.3 is calibrated so a few-mN/m tension yields a \textasciitilde{}mV-scale activation shift and a tens-of-percent Na\textsuperscript{+} current change, the order Kubanek report (not a directly tabulated value). The neuron field k\_mechano\_nav defaults to 0 (pathway off) to preserve the baseline neuron; set it to K\_MECHANO\_NAV to enable. \\
        \texttt{v\_eval} & Assumption & Near-threshold reference voltage [mV] at which nav\_current\_gain reads the open fraction. \\
        \bottomrule
    \end{tabular}
\end{table}

\begin{table}[H]
    \caption{Source classification: MechanoSynapse.}
    \label{tab:prov_mechanosynapse}
    \footnotesize
    \centering
    \begin{tabular}{@{}llp{6.5cm}@{}}
        \toprule
        Parameter & Source & Reference / notes \\
        \midrule
        \texttt{E\_ampa} & Literature & AMPA non-selective cation reversal, \textasciitilde{}0 mV \\
        \texttt{K\_ampa} & Literature & AMPA is low-affinity (glutamate EC50 \textasciitilde{} hundreds of $\mu$M) --- AMPA glutamate half-saturation [$\mu$M]. \\
        \texttt{Pr\_baseline} & Literature & Tyler 2008; central-synapse release probability range 0.1--0.3 --- Resting vesicle release probability at zero tension. \\
        \texttt{tau\_glu} & Literature & Cleft glutamate clearance / uptake time constant (few ms) --- First-order glutamate pool decay [ms]. \\
        \texttt{g\_ampa} & Calibrated & Post-synaptic AMPA conductance density [mS/cm\textsuperscript{2}]; calibrated with glu\_scale. \\
        \texttt{glu\_scale} & Calibrated & Sustained cleft glutamate [$\mu$M] at full release; calibrated with g\_ampa so the resting (baseline-Pr) state is sub-threshold and US-driven release brings the post neuron to firing. \\
        \texttt{Pr\_max} & Estimate & Tyler 2012 (mechanical perturbation raises release probability) --- Saturating release probability at high presynaptic tension. \\
        \texttt{T\_half\_pre} & Estimate & Tyler 2012 --- Half-modulation presynaptic tension [mN/m]; in the MS-channel range. \\
        \texttt{k\_pre} & Estimate & Boltzmann slope [mN/m] of the release-probability sigmoid. \\
        \bottomrule
    \end{tabular}
\end{table}

\begin{table}[H]
    \caption{Source classification: cortical\_pyramidal\_3comp.}
    \label{tab:prov_cortical_pyramidal_3comp}
    \footnotesize
    \centering
    \begin{tabular}{@{}llp{6.5cm}@{}}
        \toprule
        Parameter & Source & Reference / notes \\
        \midrule
        \texttt{dendrite\_piezo\_scale} & Literature & Lewis \& Grandl 2015; Dalghi 2019 --- Qualitative dendrite-heavy Piezo1 expression \\
        \texttt{g\_Na\_ais} & Literature & Kole 2008 (lower end of 2000-5000 mS/cm\textsuperscript{2} range) --- Default chosen for RK4 stability at dt = 1 $\mu$s; user can raise to 5000 \\
        \texttt{g\_Na\_dend} & Literature & Stuart \& Sakmann 1995 \\
        \texttt{g\_Na\_soma} & Literature & Mainen \& Sejnowski 1996 \\
        \texttt{area\_ratios} & Assumption & Pinsky-Rinzel-style lumped ratios A\_d:A\_s:A\_a = 5:1:0.05 \\
        \texttt{g\_dend\_soma} & Assumption & Lumped axial conductance, mS/cm\textsuperscript{2} of soma \\
        \texttt{g\_soma\_ais} & Assumption & Tighter coupling than dend-soma, reflecting AIS proximity \\
        \bottomrule
    \end{tabular}
\end{table}
